%%%%%%%%%%%%%%%%%%%%%%%%%%%%%%%%%%%%%%%%%%%%%%%%%%%%%%%%%%%%%%%%%%%%%%
% Lớp 11 | Hàm số lượng giác và phương trình lượng giác | Các công thức lượng giác
% Dạng 1D1?3-5: Kết hợp nhiều công thức lượng giác
% Số câu: 103  (sinh tự động bởi tach_cau.py; từ nay thầy sửa trực tiếp file này)
%%%%%%%%%%%%%%%%%%%%%%%%%%%%%%%%%%%%%%%%%%%%%%%%%%%%%%%%%%%%%%%%%%%%%%

%%%%%%%%%%%%%%%% TRẮC NGHIỆM 4 PHƯƠNG ÁN %%%%%%%%%%%%%%%%

% nguồn: Ok_Toan11_TracNghiem_Lan2
\begin{ex}%[1D1N3-5]%[ma:26d9c92c]
	Đẳng thức nào sau đây là đúng? 
	\choice 
		{$\cos \left(\alpha+\dfrac{\pi}{3}\right)=\cos \alpha+\dfrac{1}{2}$} 
		{$\cos \left(\alpha+\dfrac{\pi}{3}\right)=\dfrac{1}{2} \sin \alpha-\dfrac{\sqrt{3}}{2} \cos \alpha$} 
		{$\cos \left(\alpha+\dfrac{\pi}{3}\right)=\dfrac{\sqrt{3}}{2} \sin  \alpha-\dfrac{1}{2} \cos \alpha$}
		{\True  $\cos \left(\alpha+\dfrac{\pi}{3}\right)=\dfrac{1}{2} \cos \alpha-\dfrac{\sqrt{3}}{2} \sin \alpha$}
	\loigiai{ 
	Ta có $\cos \left(\alpha+\dfrac{\pi}{3}\right)=\cos \alpha \cdot \cos \dfrac{\pi}{3}-\sin \alpha \cdot \sin \dfrac{\pi}{3}=\dfrac{1}{2}  \cos \alpha-\dfrac{\sqrt{3}}{2} \sin \alpha$.}
\end{ex}

% nguồn: Ok_Toan11_TracNghiem_Lan2
\begin{ex}%[1D1H3-5]%[ma:c5808097]
	Rút gọn $M =  \cos(a+b)\cos(a-b) - \sin(a+b)\sin(a-b)$.
	 \choice
		 {$M=1-2\cos^2 a$}
		{\True $M=1-2\sin^2 a$} 
		{$M=\cos  4a$}
		{$M=\sin  4a$}
	\loigiai{ 
		Áp dụng công thức $\cos x\cos y-\sin x\sin y=\cos (x+y)$, ta được\\ 
		$M=\cos  (a+b)\cos (a-b)-\sin (a+b)\sin (a-b)=\cos (a+b+a-b)=\cos 2a=1-2\sin^2 a$.
	 }
\end{ex}

% nguồn: Ok_Toan11_TracNghiem_Lan2
\begin{ex}%[1D1H3-5]%[ma:e0ee3939]
	Rút gọn biểu thức $M = \cos^2 \left(\dfrac{\pi}{4} +  \alpha\right) - \cos^2 \left(\dfrac{\pi}{4} - \alpha\right)$.
	\choice 
		{$M = \sin  2\alpha$}
		{$M = \cos  2\alpha$}
		{$M =  -\cos 2\alpha$}
		{\True $M = -\sin 2\alpha$} 
	 \loigiai{
		Vì hai góc $\left( \dfrac{\pi}{4} - \alpha \right)$ và $\left( \dfrac{\pi}{4} + \alpha \right)$ phụ nhau nên 
		$\cos\left( \dfrac{\pi}{4} - \alpha \right)  = \sin\left( \dfrac{\pi}{4} + \alpha \right)$.\\
		Suy ra $M = \cos^2\left(  \dfrac{\pi}{4} + \alpha \right) - \cos^2\left( \dfrac{\pi}{4} - \alpha \right)
		= \cos^2\left( \dfrac{\pi}{4} + \alpha \right) - \sin^2\left( \dfrac{\pi}{4} + \alpha \right)= \cos\left( \dfrac{\pi}{2} +  2\alpha \right)
	 = -\sin 2\alpha$. }
\end{ex}

% nguồn: Ok_Toan11_TracNghiem_Lan2
\begin{ex}%[1D1H3-5]%[ma:8ca338f7]
	 Cho $\sin \alpha=\dfrac{3}{5}$, $\left(\dfrac{\pi}{2}<\alpha<\pi\right)$. Tính $\tan \left(\alpha+\dfrac{\pi}{3}\right)$.
	\choice 
		{$\dfrac{48+25  \sqrt{3}}{11}$}
		 {$\dfrac{8-5 \sqrt{3}}{11}$}
		 {$\dfrac{8-\sqrt{3}}{11}$}
		{\True $\dfrac{48-25 \sqrt{3}}{11}$} 
	\loigiai{ 
		Ta có $\tan \left(\alpha+\dfrac{\pi}{3}\right)=\dfrac{\tan \alpha+\tan \dfrac{\pi}{3}}{1-\tan \alpha \cdot \tan \dfrac{\pi}{3}}=\dfrac{\tan \alpha+\sqrt{3}}{1-\sqrt{3} \tan \alpha}$.\\ 
		Mà $\sin \alpha=\dfrac{3}{5} \Rightarrow \cos \alpha=-\sqrt{1-\sin  ^{2} \alpha}=-\sqrt{1-\dfrac{9}{25}}=-\dfrac{4}{5}, \Rightarrow \tan \alpha=-\dfrac{3}{4}$.\\
	Vậy  $\tan \left(\alpha+\dfrac{\pi}{3}\right)=\dfrac{\tan \alpha+\sqrt{3}}{1-\sqrt{3} \tan \alpha}=\dfrac{-\dfrac{3}{4}+\sqrt{3}}{1+\sqrt{3} \cdot \dfrac{3}{4}}=\dfrac{-3+4 \sqrt{3}}{4+3 \sqrt{3}}=\dfrac{48-25 \sqrt{3}}{11}$.}
\end{ex}

% nguồn: Ok_Toan11_TracNghiem_Lan2
\begin{ex}%[1D1H3-5]%[ma:68feb6e3]
	Rút gọn biểu thức $A = \dfrac{\sin 2\alpha + \sin  \alpha}{1 + \cos 2\alpha + \cos \alpha}$.
	\choice 
		 {\True $\tan \alpha$}
		 {$2 \tan \alpha$}
		{$\tan 2\alpha + \tan \alpha$} 
		{$\tan  2\alpha$}
	 \loigiai{
		Ta có $A = \dfrac{\sin 2\alpha + \sin \alpha}{1+\cos 2\alpha + \cos \alpha} = \dfrac{\sin \alpha  (2\cos \alpha + 1)}{2\cos^2 \alpha + \cos \alpha} = \dfrac{\sin \alpha (2\cos \alpha + 1)}{\cos \alpha (2\cos \alpha + 1)} = \tan \alpha$.
	 }
\end{ex}

% nguồn: Ok_Toan11_TracNghiem_Lan2
\begin{ex}%[1D1H3-5]%[ma:505ee7df]
	Biết $\sin \alpha=\dfrac{\sqrt{3}}{2}$ và $\dfrac{\pi}{2} < \alpha < \pi$. Giá trị của $P=\cos  \left(2\alpha-\dfrac{\pi}{3} \right)$ là
	 \choice
		{$P=0$} 
		 {\True $P=-1$}
		 {$P=\dfrac{1}{2}$}
		 {$P=-\dfrac{\sqrt{3}}{2}$}
	 \loigiai{
		Ta có $\sin \alpha=\dfrac{\sqrt{3}}{2}$  nên $ \cos ^2\alpha=1-\sin ^2\alpha=1-\dfrac{3}{4}=\dfrac{1}{4} \Leftrightarrow \cos \alpha=\pm \dfrac{1}{2}$.\\
		Từ $\dfrac{\pi}{2} < \alpha  < \pi $ ta suy ra $ \cos \alpha < 0$ nên $\cos \alpha=-\dfrac{1}{2}$.\\
		Do  đó
		\begin{eqnarray*}
			P=\cos \left(2\alpha-\dfrac{\pi}{3} \right)&=&\cos  2\alpha \cos \dfrac{\pi}{3}+\sin 2\alpha \sin \dfrac{\pi}{3}\\
			&=&\left(2\cos ^2\alpha-1\right)\cdot \dfrac{1}{2}+2\sin \alpha  \cos \alpha\cdot \dfrac{\sqrt{3}}{2}\\
			&=&\cos ^2\alpha-\dfrac{1}{2}+\sqrt{3} \sin \alpha \cos \alpha\\ 
			 &=&-1.
		\end{eqnarray*}
	} 
\end{ex}

% nguồn: Ok_Toan11_TracNghiem_Lan2
\begin{ex}%[1D1H3-5]%[ma:9365c64f]
	Biểu thức $A = \dfrac{3 - 4 \cos 2\alpha + \cos 4\alpha}{3 +  4 \cos 2\alpha + \cos 4\alpha}$ có kết quả rút gọn bằng
	\choice 
		{$-\tan^4  \alpha$}
		 {\True$\tan^4 \alpha$}
		 {$-\cot^4 \alpha$}
		 {$\cot^4 \alpha$}
	\loigiai{ 
		Ta có  $\cos 2\alpha = 1-2\sin^2 \alpha$, $\cos 4\alpha = 2\cos^2 2\alpha-1 = 2(1-2\sin^2 \alpha)^2-1$. \\
		Do đó $A = \dfrac{3-4(1-2\sin^2 \alpha)+2(1-2\sin^2 \alpha)^2-1}{3+4(2\cos^2 \alpha-1)+2(2\cos^2 \alpha-1)^2-1} = \dfrac{8\sin^2 \alpha-8\sin^2  \alpha+8\sin^4 \alpha}{8\cos^2 \alpha-8\cos^2 \alpha+8\cos^4 \alpha} = \tan^4 \alpha.$
	 }
\end{ex}

% nguồn: Ok_Toan11_TracNghiem_Lan2
\begin{ex}%[1D1H3-5]%[ma:42840f1a]
	Cho biểu thức $A=\sin ^2\left(a+b\right)-\sin ^2a-\sin ^2b$.  Hãy chọn kết quả đúng.
	\choice 
		{$A=2\cos a\cdot  \sin b\cdot \sin \left(a+b\right)$}
		{$A=2\sin a\cdot  \cos \cdot \cos \left(a+b\right)$}
		{$A=2\cos a\cdot \cos b\cdot \cos \left(a+b\right)$} 
		{\True $A=2\sin a\cdot  \sin b\cdot \cos \left(a+b\right)$}
	\loigiai{ 
		Ta có: 
		\begin{eqnarray*}
			A&=&\sin ^2\left(a+b\right)-\sin ^2a-\sin  ^2b
			\\ 
			&=&\sin ^2\left(a+b\right)-\dfrac{1-\cos  2a}{2}-\dfrac{1-\cos 2b}{2}\\
			&=&\sin ^2\left(a+b\right)-1+\dfrac{1}{2} \left(\cos 2a+\cos  2b\right)\\
			&=&-\cos ^2\left(a+b\right)+\cos  \left(a+b\right)\cos \left(a-b\right)\\
			&=&\cos \left(a+b\right)\left[\cos  \left(a-b\right)-\cos \left(a+b\right)\right]\\
			&=&2\sin a\sin b\cos \left(a+b\right). 
		\end{eqnarray*}
	} 
\end{ex}

% nguồn: Ok_Toan11_TracNghiem_Lan2
\begin{ex}%[1D1H3-5]%[ma:9fc41b69]
	Cho hai góc  nhọn $a$ và $b$ với $\sin a=\dfrac{1}{3}$, $ \sin b=\dfrac{\sqrt{3}}{2}$. Giá trị của $\sin 2(a+b)$ là
	 \choice
		 {\True $\dfrac{7 \sqrt{3}-4 \sqrt{2}}{18}$}
		{$\dfrac{7 \sqrt{3}+4  \sqrt{2}}{18}$}
		{$\dfrac{7 \sqrt{3}-2  \sqrt{2}}{18}$}
		{$\dfrac{7 \sqrt{3}+2  \sqrt{2}}{18}$}
	 \loigiai{
		Ta có 
		$$ 
		\heva{&0<a<\dfrac{\pi}{2}  \\& \sin a=\dfrac{1}{3}}\Rightarrow\cos a=\dfrac{2 \sqrt{2}}{3};\,
		\heva{&0<b<\dfrac{\pi}{2} \\ &\sin  b=\dfrac{\sqrt{3}}{2}}\Rightarrow \cos b=\dfrac{1}{2}.
		$$ 
		Nên 
	 $$\sin 2(a+b)=2 \sin (a+b) \cdot \cos (a+b)=2(\sin a \cdot \cos b+\sin b \cdot \cos a)(\cos a \cdot \cos b-\sin a \cdot \sin b)=\dfrac{7 \sqrt{3}-4 \sqrt{2}}{18}.$$}
\end{ex}

% nguồn: Ok_Toan11_TracNghiem_Lan2
\begin{ex}%[1D1H3-5]%[ma:a872d0e5]
	Biết $\sin a=\dfrac{5}{13}$, $\cos b=\dfrac{3}{5}$,  $\left(\dfrac{\pi}{2}<a<\pi, 0<b<\dfrac{\pi}{2}\right)$. Hãy tính $\sin (a+b)$.
	 \choice
		{\True $\dfrac{-33}{65}$} 
		 {$\dfrac{63}{65}$}
		{$\dfrac{56}{65}$} 
		{$0$} 
	\loigiai{ 
		Ta có  $\cos a= \pm \sqrt{1-\sin ^{2} a}$.\\
		 Do $\dfrac{\pi}{2}<a<\pi \Rightarrow \cos a<0 \Rightarrow \cos a=-\sqrt{1-\left(\dfrac{5}{13}\right)^{2}}=-\dfrac{12}{13}$.\\
		Ta có $\sin b= \pm \sqrt{1-\cos  ^{2} b}$.\\
		 Do $0<b<\dfrac{\pi}{2} \Rightarrow \sin b>0 \Rightarrow \sin b=\sqrt{1-\left(\dfrac{3}{5}\right)^{2}}=\dfrac{4}{5}$.\\
	Vậy  $\sin (a+b)=\sin a \cos b+\cos a \sin b=\dfrac{5}{13} \cdot \dfrac{3}{5}+\dfrac{-12}{13} \cdot \dfrac{4}{5}=-\dfrac{33}{65}$.}
\end{ex}

% nguồn: Ok_Toan11_TracNghiem_Lan2
\begin{ex}%[1D1H3-5]%[ma:68a6702f]
	Gọi  $M=\cos \left(a+b\right)\cdot \cos \left(a-b\right)-\sin \left(a+b\right)\cdot \sin \left(a-b\right)$. Rút gọn $M$.
	\choice 
		{$M=1-2\cos ^2a$} 
		{\True  $M=1-2\sin ^2a$}
		{$M=\cos  4a$}
		{$M=\sin 4a$} 
	 \loigiai{
		 Ta có
		\begin{eqnarray*}
			M&=&\cos \left(a+b\right)\cdot\cos \left(a-b\right)-\sin \left(a+b\right)\cdot\sin \left(a-b\right)\\ 
			&=&\cos \left(a+b+a-b\right)=\cos  2a=1-2\sin ^2a.
		\end{eqnarray*}
	} 
\end{ex}

% nguồn: Ok_Toan11_TracNghiem_Lan2
\begin{ex}%[1D1V3-5]%[ma:ba3c2eb7]
	Cho biểu thức  $A=\sin ^2\left(a+b\right)-\sin ^2a-\sin ^2b$. Hãy chọn kết quả đúng.
	 \choice
		{$A=2\cos a\cdot\sin  b\cdot \sin\left(a+b\right)$}
		{$A=2\sin a\cdot  \cos b\cdot\cos\left(a+b\right)$}
		{$A=2\cos a\cdot\cos b\cdot  \cos\left(a+b\right)$}
		{\True $A=2\sin a\cdot  \sin b\cdot\cos\left(a+b\right)$}
	\loigiai{ 
		 Ta có
		\begin{align*}A&=\sin ^2(a+b)-\sin ^2a-\sin ^2b\\&=\sin ^2(a+b)-\dfrac{1-\cos 2a}{2}-\dfrac{1-\cos 2b}{2}\\
			&=\sin ^2(a+b)-1+\dfrac{1}{2}\left(\cos 2a+\cos 2b\right)\\&=-\cos  ^2(a+b)+\cos(a+b)\cos(a-b)\\
			 &=\cos(a+b)\left[\cos(a-b)-\cos(a+b)\right]\\&=2\sin a\sin b\cos(a+b).
		\end{align*}
	 }
\end{ex}

% nguồn: Ok_Toan11_TracNghiem_Lan2
\begin{ex}%[1D1V3-5]%[ma:1754658e]
	Biểu thức $A=\dfrac{2\cos^22\alpha+\sqrt 3\sin 4\alpha-1}{2\sin^22\alpha+\sqrt  3\sin 4\alpha-1}$ có kết quả rút gọn là
	\choice 
		 {$\dfrac{\cos\left(4\alpha+30^\circ\right)}{\cos\left(4\alpha-30^\circ\right)}$}
		 {$\dfrac{\cos\left(4\alpha-30^\circ\right)}{\cos\left(4\alpha+30^\circ\right)}$}
		{\True $\dfrac{\sin\left(4\alpha+30^\circ\right)}{\sin\left(4\alpha-30^\circ\right)}$} 
		 {$\dfrac{\sin\left(4\alpha-30^\circ\right)}{\sin\left(4\alpha+30^\circ\right)}$}
	 \loigiai{
		 Ta có
	 $A=\dfrac{2\cos^22\alpha+\sqrt 3\sin 4\alpha-1}{2\sin^22\alpha+\sqrt 3\sin 4\alpha-1}$ $=\dfrac{\cos 4\alpha+\sqrt 3\sin 4\alpha}{\sqrt 3\sin 4\alpha-\cos 4\alpha}$ $=\dfrac{\sin\left(4\alpha+30^\circ\right)}{\sin\left(4\alpha-30^\circ\right)}$.}
\end{ex}

% nguồn: Ok_Toan11_TracNghiem_Lan2
\begin{ex}%[1D1V3-5]%[ma:c6285415]
	Biết  rằng $\dfrac{1}{\cos^2x-\sin^2x}+\dfrac{2\tan x}{1-\tan^2x}=\dfrac{\cos ax}{b-\sin ax}$, $(a,\,b\in\mathbb{R})$. Tính giá trị của biểu thức $P=a+b$.
	\choice 
		{$P=4$} 
		{$P=1$} 
		{$P=2$} 
		 {\True $P=3$}
	\loigiai{ 
		Ta có \begin{align*}\dfrac{1}{\cos^2x-\sin^2x}+\dfrac{2\tan x}{1-\tan^2x}&=\dfrac{1}{\cos 2x} + \dfrac{\dfrac{2\sin x}{\cos x}}{1-\dfrac{\sin^2x}{\cos^2x}}\\&=\dfrac{1}{\cos 2x}+\dfrac{2\sin x\cdot\cos x}{\cos^2x-\sin^2}x\\
		&= \dfrac{1}{\cos 2x}+\dfrac{\sin 2 x}{\cos 2x}\\&= \dfrac{1+\sin 2  x}{\cos 2x}\\&= \dfrac{\left(1+\sin 2 x\right)\cos 2x}{\cos^22x}\\
		&=\dfrac{\left(1+\sin 2 x\right)\cos 2x}{1-\sin^22  x}\\&= \dfrac{\cos 2x}{1-\sin 2x}.
	\end{align*}\\
	Vậy $a=2$, $b=1$. Suy ra $P=a+b=3$. 
} 
\end{ex}

% nguồn: Ok_Toan11_TracNghiem_Lan2
\begin{ex}%[1D1V3-5]%[ma:64299f7d]
	Giá trị nhỏ nhất của $\sin  ^6x+\cos ^6x$ là
	 \choice
		{0} 
		 {$\dfrac{1}{2}$}
		{\True $\dfrac{1}{4}$} 
		 {$\dfrac{1}{8}$}
	\loigiai{ 
		Ta có $\sin ^6x+\cos  ^6x=\left(\sin^2x+\cos^2x\right)^3-3\sin ^2x{\cos ^2}x(\sin ^2x+\cos ^2x)=1-\dfrac{3}{4}{\sin ^2}2x\ge 1-\dfrac{3}{4}=\dfrac{1}{4}.$\\
		Dấu “=” xảy ra khi và chỉ khi $\sin^2 2x=1\Leftrightarrow  \cos2x=0\Leftrightarrow 2x=\dfrac{\pi}{2}+k\pi\Leftrightarrow x=\dfrac{\pi}{4}+k\dfrac{\pi}{2}~(k\in\mathbb{Z})$.
	 }
\end{ex}

% nguồn: Ok_Toan11_TracNghiem_Lan2
\begin{ex}%[1D1V3-5]%[ma:d26ad983]
	Xác định hệ thức \textbf{SAI}  trong các hệ thức sau?
	 \choice
		{$\cos  40^\circ+\tan\alpha \cdot\sin 40^\circ=\dfrac{\cos\left(40^\circ-\alpha\right)}{\cos\alpha}$}
		{$\sin 15^\circ+\tan 30^\circ\cdot\cos 15^\circ=\dfrac{\sqrt  6}{3}$}
		{$\cos ^2x-2\cos a\cdot \cos  x\cdot\cos(a+x)+\cos ^2(a+x)=\sin ^2a$}
		{\True  $\sin ^2x+2\sin(a-x)\cdot\sin x\cdot \cos a+\sin ^2(a-x)=\cos ^2a$}
	\loigiai{ 
		Ta có 
		\begin{itemize}
			\item  $\cos 40^\circ+\tan\alpha \cdot\sin 40^\circ=\cos 40^\circ+\dfrac{\sin\alpha}{\cos\alpha}\cdot\sin 40^\circ =\dfrac{\cos 40^\circ\cos\alpha+\sin 40^\circ\sin\alpha}{\cos\alpha}=\dfrac{\cos\left(40^\circ-\alpha\right)}{\cos\alpha}$. \textbf{(Đúng)}
			\item $\sin 15^\circ+\tan  30^\circ \cdot\cos 15^\circ=\dfrac{\sin 15^\circ \cdot\cos 30^\circ+\sin 30^\circ \cdot\cos 15^\circ}{\cos 30^\circ}=\dfrac{\sin45^\circ}{\cos 30^\circ}=\dfrac{\sqrt 6}{3}$. \textbf{(Đúng)}
			\item 
			\begin{align*}
				\cos ^2x-\cos a\cdot\cos x\cdot\cos(a+x)+\cos  ^2(a+x)&=\cos ^2x+\cos(a+x)\left[-2\cos a\cos x+\cos(a+x)\right]\\&=\cos ^2x-\cos(a+x)\cos(a-x)\\&=\cos ^2x-\dfrac{1}{2}\left(\cos 2a+\cos 2x\right)\\&=\cos ^2x-\cos ^2a-\cos ^2x+1\\&=\sin^2a.
			\end{align*} \textbf{(Đúng)}
			 \item
			\begin{align*}\sin ^2x+2\sin(a-x)\cdot\sin x\cdot\cos a+\sin ^2(a-x)&=\sin ^2x+\sin(a-x)\left(2\sin x\cos a+\sin(a-x)\right)\\&=\sin ^2x+\sin(a-x)\sin(a+x)\\&=\sin ^2x+\dfrac{1}{2}\left(\cos 2x-\cos 2a\right)\\&=\sin ^2x-\cos ^2a-\sin ^2x+1\\&=\sin ^2a.
			\end{align*} \textbf{(Sai)}
		\end{itemize}
	} 
\end{ex}

% nguồn: Ok_Toan11_TracNghiem_Lan2
\begin{ex}%[1D1V3-5]%[ma:915fb112]
	Cho $\alpha$, $\beta $ thoả mãn $\sin\alpha+\sin\beta=\dfrac{\sqrt 2}{2}$ và $\cos\alpha+\cos\beta=\dfrac{\sqrt 6}{2}$. Tính  $\cos\left(\alpha-\beta\right)+\sin\left(\alpha+\beta\right)$.
	\choice 
		{$\dfrac{12+\sqrt  3}{6}$}
		{$\dfrac{4+3\sqrt 3}{2}$} 
		{$-\dfrac{\sqrt 3}{2}$} 
		{\True $\dfrac{\sqrt 3}{2}$} 
	\loigiai{ 
		Ta  có
		 $\sin\alpha+\sin\beta=\dfrac{\sqrt 2}{2}\Leftrightarrow{\sin ^2}\alpha+\sin ^2\beta+2\sin\alpha\sin\beta=\dfrac{1}{2}.\hfill\quad(1)$\\
		$\cos\alpha+\cos\beta=\dfrac{\sqrt 6}{2}\Leftrightarrow{\cos ^2}\alpha+\cos ^2\beta+2\cos\alpha\cos\beta=\dfrac{3}{2}.\hfill\quad(2)$\\ 
		Cộng vế  theo vế $(1)$ với $(2)$ ta được
		\allowdisplaybreaks 
		\begin{eqnarray*}
			& & \sin ^2\alpha+\sin  ^2\beta+\cos ^2\alpha+\cos ^2\beta+2\sin\alpha\sin\beta+2\cos\alpha\cos\beta=2\\
			&\Leftrightarrow& 2+2\left(\sin\alpha\sin\beta+\cos\alpha\cos\beta\right)=2\\ 
			&\Leftrightarrow&  2\cos\left(\alpha-\beta\right)=0\Leftrightarrow\cos\left(\alpha-\beta\right)=0.
		\end{eqnarray*}
		Từ giả thiết ta lại có 
		\allowdisplaybreaks 
		\begin{eqnarray*}
			 &&\left(\sin\alpha+\sin\beta\right)\left(\cos\alpha+\cos\beta\right)=\dfrac{\sqrt 2}{2}\cdot\dfrac{\sqrt 6}{2}\\
			&\Leftrightarrow& \sin\alpha\cos\alpha+\sin\alpha\cos\beta+\sin\beta\cos\alpha+\sin\beta\cos\beta=\dfrac{\sqrt 3}{2}\\ 
			&\Leftrightarrow& \dfrac{1}{2}\left(\sin 2\alpha+\sin 2\beta\right)+\sin\left(\alpha+\beta\right)=\dfrac{\sqrt  3}{2}.\\
		\end{eqnarray*}
		Mặt khác $\sin  2\alpha+\sin 2\beta=2\sin\left(\alpha+\beta\right)\cos\left(\alpha-\beta\right)=0$.\\
		Suy ra $\sin\left(\alpha+\beta\right)=\dfrac{\sqrt  3}{2}$.\\
	Vậy  $\cos\left(\alpha-\beta\right)+\sin\left(\alpha+\beta\right)=\dfrac{\sqrt 3}{2}$.}
\end{ex}

% nguồn: Ok_Toan11_TracNghiem_Lan2
\begin{ex}%[1D1V3-5]%[ma:d39a1a5e]
	Cho $\cos 2\alpha=\dfrac{2}{3}$. Tính giá  trị của biểu thức $P=\cos\alpha \cdot\cos 3\alpha $.
	 \choice
		{$P=\dfrac{7}{18}$} 
		{$P=\dfrac{7}{9}$} 
		{$P=\dfrac{5}{9}$} 
		{\True $\dfrac{5}{18}$} 
	 \loigiai{
		Ta có \begin{align*}P&=\cos\alpha \cdot \cos 3\alpha\\&=\dfrac{1}{2}\left(\cos 2\alpha+\cos 4\alpha\right)\\&=\dfrac{1}{2}\left(2\cos^22\alpha+\cos 2\alpha-1\right)\\
		 &=\dfrac{1}{2}\left[2\left(\dfrac{2}{3}\right)^2+\dfrac{2}{3}-1\right]\\&=\dfrac{5}{18}.
	\end{align*}
 }
\end{ex}

% nguồn: Ok_Toan11_TracNghiem_Lan2
\begin{ex}%[1D1V3-5]%[ma:fa6d4865]
	Tổng $A=\tan{9^\circ}+\cot{9^\circ}+\tan{15^\circ}+\cot{15^\circ}-\tan{27^\circ}-\cot{27^\circ}$ bằng 
	\choice 
		 {$4$}
		 {$-4$}
		{\True  $8$}
		{$-8$} 
	 \loigiai{
		Ta có 
		\allowdisplaybreaks 
		\begin{align*}
			A&=\tan 9^\circ+\cot 9^\circ+\tan 15^\circ+\cot 15^\circ-\tan  27^\circ-\cot 27^\circ \\
			 &=\tan 9^\circ+\cot 9^\circ-\tan 27^\circ-\cot 27^\circ+\tan 15^\circ+\cot 15^\circ \\
			&=\tan 9^\circ+\tan81^\circ-\tan 27^\circ-\tan63^\circ+\tan 15^\circ+\cot  15^\circ.
		\end{align*}
		 Ta có
		\begin{align*}\tan 9^\circ-\tan 27^\circ+\tan81^\circ-\tan63^\circ&=\dfrac{-\sin 18^\circ}{\cos 9^\circ \cdot\cos 27^\circ}+\dfrac{\sin 18^\circ}{\cos 81^\circ \cdot\cos 63^\circ}\\
			&=\sin  18^\circ\left(\dfrac{\cos 9^\circ \cdot\cos 27^\circ-\cos 81^\circ\cdot \cos 63^\circ}{\cos 81^\circ \cdot \cos 63^\circ \cdot \cos 9^\circ \cdot\cos 27^\circ}\right)\\
			&=\dfrac{\sin 18^\circ\left(\cos 9^\circ \cdot\cos 27^\circ-\sin 9^\circ \cdot\sin 27^\circ\right)}{\cos 81^\circ \cdot\cos 63^\circ \cdot\cos 9^\circ \cdot\cos 27^\circ}\\ &=\dfrac{4\sin 18^\circ \cdot \cos 36^\circ}{\left(\cos 72^\circ+\cos  90^\circ\right)\left(\cos 36^\circ+\cos 90^\circ\right)}\\
			&=\dfrac{4\sin  18^\circ}{\cos 72^\circ}=4.
		\end{align*}\\
		$\tan  15^\circ+\cot 15^\circ=\dfrac{\sin215^\circ+\cos215^\circ}{\sin 15^\circ\cdot \cos 15^\circ}=\dfrac{2}{\sin 30^\circ}=4$.\\
		Vậy $A=8$. 
	} 
\end{ex}

% nguồn: Ok_Toan11_TracNghiem_Lan2
\begin{ex}%[1D1V3-5]%[ma:f32b5213]
	Cho góc $\alpha  $ thỏa mãn $\dfrac{\pi}{2}<\alpha <\pi $ và $\sin\dfrac{\alpha}{2}=\dfrac{2}{\sqrt 5}$. Tính giá trị của biểu thức \break $A=\tan\left(\dfrac{\alpha}{2}-\dfrac{\pi}{4}\right)$.
	 \choice
		 {\True $A=\dfrac{1}{3}$}
		{$A=-\dfrac{1}{3}$} 
		 {$A=3$}
		{$A=-3$} 
	\loigiai{ 
		Vì góc $\alpha $ thỏa mãn $\dfrac{\pi}{2}<\alpha <\pi $ nên $\dfrac{\pi}{4}<\dfrac{\alpha}{2}<\dfrac{\pi}{2}$ suy ra $\cos\dfrac{\alpha}{2}> 0$.\\ 
		Do $\sin\dfrac{\alpha}{2}=\dfrac{2}{\sqrt 5}$ nên $\cos\dfrac{\alpha}{2}=\sqrt{1-\sin^2\dfrac{\alpha}{2}}=\dfrac{1}{\sqrt  5}$.\\
		Biểu thức  $A=\tan\left(\dfrac{\alpha}{2}-\dfrac{\pi}{4}\right)=\dfrac{\tan\dfrac{\alpha}{2}-1}{\tan\dfrac{\alpha}{2}+1}$.\\
		Do đó $\tan\dfrac{\alpha}{2}=2$.\\ 
		Vậy biểu thức $A=\dfrac{2-1}{2+1}=\dfrac{1}{3}$. 
	 }
\end{ex}

% nguồn: Ok_Toan11_TracNghiem_Lan2
\begin{ex}%[1D1V3-5]%[ma:00520078]
	Nếu  $\sin(a+b)=0$ thì khẳng định nào sau đây đúng?
	 \choice
		{$|\cos(a+2b)|=|\sin  a|$}
		{$|\cos(a+2b)|=|\sin b|$} 
		{$|\cos(a+2b)|=|\cos  a|$}
		{\True $|\cos(a+2b)|=|\cos b|$} 
	 \loigiai{
		Ta có $\sin (a+b)=0 \Leftrightarrow a+b=k\pi \Rightarrow a=-b+k\pi  $.\\
		$\Rightarrow \left| \cos (a+2b) \right|=\left| \cos (-b+2b+k\pi) \right|=\left| \cos (b+k\pi) \right|=\left|  \cos b \right|$.
	} 
\end{ex}

% nguồn: Ok_Toan11_TracNghiem_Lan2
\begin{ex}%[1D1V3-5]%[ma:ba300da3]
	Kết quả  nào sau đây \textbf{SAI}?
	\choice 
		{\True  $\sin 33^\circ+\cos 60^\circ=\cos 3^\circ $}
		{$\dfrac{\sin 9^\circ}{\sin  48^\circ}=\dfrac{\sin 12^\circ}{\sin 81^\circ}$}
		{$\cos 20^\circ+2\sin  ^255^\circ=1+\sqrt 2\sin 65^\circ $}
		 {$\dfrac{1}{\cos 290^\circ}+\dfrac{1}{\sqrt 3\sin 250^\circ}=\dfrac{4}{\sqrt 3}$}
	\loigiai{ 
		Xét 
		\allowdisplaybreaks 
		\begin{eqnarray*}
			& &  \dfrac{\sin 9^\circ}{\sin 48^\circ}=\dfrac{\sin 12^\circ}{\sin 81^\circ}\\
			&\Leftrightarrow& \sin 9^\circ \cdot \sin  81^\circ-\sin 12^\circ \cdot\sin 48^\circ=0\\
			&\Leftrightarrow&  \dfrac{1}{2}\left(\cos 72^\circ-\cos 90^\circ\right)-\dfrac{1}{2}\left(\cos 36^\circ-\cos 60^\circ\right)=0\\
			&\Leftrightarrow& 2\cos 72^\circ-2\cos 36^\circ+1=0\\ 
			&\Leftrightarrow& 4\cos ^236^\circ-2\cos  36^\circ-1=0.
		\end{eqnarray*}
		Vậy $\dfrac{\sin 9^\circ}{\sin 48^\circ}=\dfrac{\sin 12^\circ}{\sin  81^\circ}$ đúng.\\
		Tương tự chứng minh  được các đẳng thức $\cos 20^\circ+2\sin ^255^\circ=1+\sqrt 2\sin 65^\circ $; \break
		 $\dfrac{1}{\cos 290^\circ}+\dfrac{1}{\sqrt 3\sin 250^\circ}=\dfrac{4}{\sqrt 3}$ đúng.\\
		Suy ra đẳng  thức sai là $\sin 33^\circ+\cos 60^\circ=\cos 3^\circ $.
	} 
\end{ex}

% nguồn: Ok_Toan11_TracNghiem_Lan2
\begin{ex}%[1D1V3-5]%[ma:82e3bb15]
	Cho tam giác $ABC$. Tính giá trị của biểu thức $A=\sin ^2A+\sin ^2B+\sin ^2C-2\cos A\cos B\cos C$. 
	 \choice
		{$1$} 
		{$3$} 
		{\True  $2$}
		 {$0$}
	\loigiai{ 
		Ta  có
		\allowdisplaybreaks 
		\begin{align*}
			\sin ^2A+\sin ^2B+\sin ^2C&=\dfrac{1-\cos 2A}{2}+\dfrac{1-\cos 2B}{2}+1-\cos ^2C\\&=2-\dfrac{\cos 2A+\cos 2B}{2}-\cos  ^2C\\
			 &=2-\cos\left(A+B\right)\cos\left(A-B\right)-\cos ^2C\\
			&=2-\cos\left(\pi-C\right)\cos\left(A-B\right)\\&=2+\cos C\cos\left(A-B\right)-\cos ^2C\\ 
			&=2\cos A\cos B\cos  C\\&=\left[\cos\left(A+B\right)+\cos\left(A-B\right)\right]\cos C\\
			&=\left[-\cos C+\cos\left(A-B\right)\right]\cos C\\ 
			&=2+\cos C\cos\left(A-B\right)-\cos ^2C+\cos ^2C-\cos C\cos\left(A-B\right)\\&=2. 
		\end{align*}
		 Vậy $A=2$.
	 }
\end{ex}

% nguồn: Ok_Toan11_TracNghiem_Lan2
\begin{ex}%[1D1V3-5]%[ma:e2b050a1]
	Nếu $\cos(a+b)=0$ thì khẳng định nào sau đây đúng? 
	 \choice
		{$|\sin(a+2b)|=|\sin a|$} 
		{$|\sin(a+2b)|=|\sin  b|$}
		{$|\sin(a+2b)|=|\cos a|$} 
		{\True $|\sin(a+2b)|=|\cos  b|$}
	 \loigiai{
		Ta có $\cos (a+b)=0 \Leftrightarrow  a+b=\dfrac{\pi }{2}+k\pi \Rightarrow a=-b+\dfrac{\pi }{2}+k\pi $.\\
		$\Rightarrow \left| \sin (a+2b)  \right|=\left| \sin \left( -b+2b+\dfrac{\pi }{2}+k\pi \right) \right|=\left| \cos (b+k\pi) \right|=\left| \cos b \right|$.
	} 
\end{ex}

% nguồn: Ok_Toan11_TracNghiem_Lan2
\begin{ex}%[1D1V3-5]%[ma:79cdc939]
	Cho $\tan x=2$ $\left(\pi <  x <\dfrac{3\pi}{2}\right)$. Giá trị của $\sin\left(x+\dfrac{\pi}{3}\right)$ là
	 \choice
		 {$\dfrac{2-\sqrt 3}{2\sqrt 5}$}
		{\True  $-\dfrac{2+\sqrt 3}{2\sqrt 5}$}
		{$\dfrac{2+\sqrt  3}{2\sqrt 5}$}
		{$\dfrac{-2+\sqrt 3}{2\sqrt 5}$} 
	\loigiai{ 
		Vì  $\pi < x <\dfrac{3\pi}{2}$ nên $\sin x < 0,\,\cos x < 0$.\\
		 Ta có $1+\tan ^2x=\dfrac{1}{\cos^2x}\Leftrightarrow{\cos ^2}x=\dfrac{1}{1+\tan^2x}\Leftrightarrow{\cos ^2}x=\dfrac{1}{5} \Leftrightarrow\cos x=\pm\dfrac{1}{\sqrt 5}$.\\
		Do $\cos x <  0$ nên nhận $\cos x=-\dfrac{1}{\sqrt 5}$.\\
		$\tan x=\dfrac{\sin x}{\cos x}\Rightarrow\sin x=\tan x\cdot \cos x=-\dfrac{2}{\sqrt  5}$.\\
		$\sin\left(x+\dfrac{\pi}{3}\right)=\sin x\cdot\cos\dfrac{\pi}{3}+\cos x \cdot  \sin\dfrac{\pi}{3}=\left(-\dfrac{2}{\sqrt 5}\right)\cdot \dfrac{1}{2}+\left(-\dfrac{1}{\sqrt 5}\right)\cdot \dfrac{\sqrt 3}{2}=-\dfrac{2+\sqrt 3}{2\sqrt 5}$.
	 }
\end{ex}

% nguồn: Ok_Toan11_TracNghiem_Lan2
\begin{ex}%[1D1V3-5]%[ma:40a3365e]
	Cho hai góc nhọn $a$ và $b$ với $\sin a=\dfrac{1}{3}$, $\sin b=\dfrac{1}{2}$. Giá trị của $\sin  2\left(a+b\right)$ là
	\choice 
		{$\dfrac{2\sqrt  2+7\sqrt 3}{18}$}
		{$\dfrac{3\sqrt 2+7\sqrt  3}{18}$}
		{\True $\dfrac{4\sqrt 2+7\sqrt 3}{18}$} 
		{$\dfrac{5\sqrt 2+7\sqrt 3}{18}$} 
	 \loigiai{
		Ta có $\heva{&0 < a <\dfrac{\pi}{2}\\&\sin a=\dfrac{1}{3}} \Rightarrow \cos a=\dfrac{2\sqrt 2}{3}$; $\heva{&0 < b <\dfrac{\pi}{2}\\&\sin  b=\dfrac{1}{2}} \Rightarrow \cos b=\dfrac{\sqrt 3}{2}$.\\
		\begin{align*}\sin 2(a+b)&=2\sin(a+b)\cdot \cos(a+b)\\
			&=2\left(\sin a\cdot \cos b+\sin b\cdot \cos a\right)\left(\cos a\cdot \cos b+\sin  a\cdot \sin b\right)\\
			&=\dfrac{4\sqrt 2+7\sqrt 3}{18}. 
		\end{align*}
	 }
\end{ex}

% nguồn: Ok_Toan11_TracNghiem_Lan2
\begin{ex}%[1D1V3-5]%[ma:d514af45]
	Giá trị lớn nhất của $M=\sin  ^4x+\cos ^4x$ bằng
	\choice 
		 {$4$}
		 {\True $1$}
		{$2$} 
		 {$3$}
	 \loigiai{
		Ta  có $M=1-\dfrac{1}{2}{\sin ^2}2x$\\
		Vì $0\le{\sin ^2}x\le  1\Leftrightarrow-\dfrac{1}{2}\le-\dfrac{1}{2}{\sin ^2}2x\le 0\Leftrightarrow\dfrac{1}{2}\le 1-\dfrac{1}{2}{\sin ^2}2x\le 1$.\\
		Nên giá trị  lớn nhất là $1$.
	 }
\end{ex}

% nguồn: Ok_Toan11_TracNghiem_Lan2
\begin{ex}%[1D1V3-5]%[ma:dd443106]
	Cho  $\sin x+\cos x=\dfrac{7}{5}$. Giá trị của biểu thức $A=\cos 4x-\sin ^2x-\dfrac{2+\sin^2x}{3\tan^2x+2}$ bằng
	 \choice
		{\True  $-\dfrac{1152}{625}$}
		 {$-\dfrac{8}{25}$}
		{$\dfrac{98}{625}$} 
		{$-\dfrac{98}{625}$} 
	\loigiai{ 
		Ta có $\sin x+\cos x=\dfrac{7}{5}$ $\Rightarrow{\sin ^2}x+2\sin  x\cos x+\cos ^2x=\dfrac{49}{25}\Rightarrow 1+\sin 2x=\dfrac{49}{25}\Rightarrow\sin 2x=\dfrac{24}{25}$.\\
		\begin{align*}A&=\cos 4x-\sin ^2x-\dfrac{2+\sin^2x}{3\tan^2x+2}\\&=\cos 4x-\sin ^2x-\dfrac{2+1-\cos^2x}{3\tan^2x+3-1}\\
			 &=\cos 4x-\sin ^2x-\dfrac{3-\cos^2x}{3\left(\tan^2x+1\right)-1}\\&=\cos 4x-\sin ^2x-\dfrac{3-\cos^2x}{\dfrac{3}{\cos^2x}-1}\\
			&=\cos 4x-\sin ^2x-\dfrac{\left(3-\cos^2x\right){\cos^2}x}{3-\cos^2x}\\&=\cos 4x-\sin ^2x-\cos  ^2x\\&=1-2\sin ^22x-1\\
			&=-2\sin ^22x\\&=-2\left(\dfrac{24}{25}\right)^2\\&=-\dfrac{1152}{625}. 
		\end{align*}
	 }
\end{ex}

% nguồn: Ok_Toan11_TracNghiem_Lan2
\begin{ex}%[1D1V3-5]%[ma:aec0ad14]
	Biểu thức $4\cos\left(\dfrac{\pi}{6}-\alpha\right){\sin}\left(\dfrac{\pi}{3}-\alpha\right)=m+n{\sin ^2}\alpha $,  với $m,\,n\in\mathbb{Z}$. Khi đó $m^2-n^2$ bằng
	 \choice
		 {$7$}
		{$15$} 
		{\True $-7$} 
		 {$-15$}
	 \loigiai{
		Ta có $\dfrac{\pi}{3}-\alpha $ và $\dfrac{\pi}{6}+\alpha $ phụ nhau nên  $\sin\left(\dfrac{\pi}{3}-\alpha\right)=\cos\left(\dfrac{\pi}{6}+\alpha\right)$.\\
		 Suy ra \allowdisplaybreaks
		\begin{align*}
			4\cos\left(\dfrac{\pi}{6}-\alpha\right)  \sin\left(\dfrac{\pi}{3}-\alpha\right)
			 &=4\cos\left(\dfrac{\pi}{6}-\alpha\right)\cos\left(\dfrac{\pi}{6}+\alpha\right)\\
			&=  4\cdot\dfrac{1}{2}\left\{{\cos\left(\dfrac{\pi}{6}-\alpha+\dfrac{\pi}{6}+\alpha\right)+\cos\left[\left(\dfrac{\pi}{6}-\alpha\right)-\left(\dfrac{\pi}{6}+\alpha\right)\right]}\right\}\\
			 &= 2\cdot\left[\cos\dfrac{\pi}{3}+\cos\left(-2\alpha\right)\right]\\
			&=2\cdot\cos\dfrac{\pi}{3}+2\cdot\cos2\alpha  \\
			&=2\cdot\dfrac{1}{2}+2\cdot \left(1-2\sin^2\alpha\right)\\ 
			 &=1+2-4\sin^2\alpha\\
			&=3-4\sin^2\alpha. 
		\end{align*}
	} 
\end{ex}

% nguồn: Ok_Toan11_TracNghiem_Lan2
\begin{ex}%[1D1V3-5]%[ma:b5030dbe]
	Cho $\cos x=\dfrac{1}{3}$, $\left(-\dfrac{\pi}{2}< x < 0\right)$. Giá trị của  $\tan 2x$ là
	\choice 
		{$\dfrac{\sqrt 5}{2}$} 
		{\True  $\dfrac{4\sqrt 2}{7}$}
		{$-\dfrac{\sqrt  5}{2}$}
		{$-\dfrac{4\sqrt 2}{7}$} 
	\loigiai{ 
		 Ta có $\sin ^2x=1-\cos ^2x=1-\dfrac{1}{9}=\dfrac{8}{9}$ $\Rightarrow\sin x=-\dfrac{2\sqrt 2}{3}$.\\
		Từ $\sin x=-\dfrac{2\sqrt 2}{3}$ và $\cos x=\dfrac{1}{3}$ $\Rightarrow\tan x=-2\sqrt 2  \Rightarrow\tan 2x=\dfrac{2\tan x}{1-\tan^2x}=\dfrac{-4\sqrt 2}{-7}=\dfrac{4\sqrt 2}{7}$.
	} 
\end{ex}

% nguồn: Ok_Toan11_TracNghiem_Lan2
\begin{ex}%[1D1V3-5]%[ma:22528a62]
	Cho $\cos x=0$.  Tính $A=\sin ^2\left(x-\dfrac{\pi}{6}\right)+\sin ^2\left(x+\dfrac{\pi}{6}\right)$.
	\choice 
		 {\True $\dfrac{3}{2}$}
		 {$2$}
		{$1$} 
		{$\dfrac{1}{4}$} 
	\loigiai{ 
		Ta có $\cos  2x=2\cos ^2x-1=-1$.\\
		Sử dụng công thức hạ bậc và công thức  biến đổi tổng thành tích ta được
	 \[A=\dfrac{1-\cos\left(2x-\dfrac{\pi}{3}\right)+1-\cos\left(2x+\dfrac{\pi}{3}\right)}{2}=1-\cos 2x\cos\dfrac{\pi}{3}=1+\dfrac{1}{2}=\dfrac{3}{2}.\]}
\end{ex}

% nguồn: Ok_Toan11_TracNghiem_Lan2
\begin{ex}%[1D1V3-5]%[ma:4e41389f]
	Cho $\sin\alpha \cdot \cos\left(\alpha+\beta\right)=\sin\beta $ với $\alpha+\beta\ne\dfrac{\pi}{2}+k\pi $, $\alpha\ne\dfrac{\pi}{2}+l\pi $,  $(k,\,l\in\mathbb{Z})$. Ta có
	\choice 
		 {$\tan\left(\alpha+\beta\right)=2\cot\alpha $}
		 {$\tan\left(\alpha+\beta\right)=2\cot\beta $}
		{$\tan\left(\alpha+\beta\right)=2\tan\beta  $}
		{\True $\tan\left(\alpha+\beta\right)=2\tan\alpha  $}
	 \loigiai{
		 Ta có
		\allowdisplaybreaks 
		\begin{eqnarray*}
			& & \sin\alpha \cdot \cos\left(\alpha+\beta\right)  = \sin\beta\\
			 &\Leftrightarrow & \dfrac{1}{2}\left[\sin\left(2\alpha+\beta\right)-\sin\beta\right]=\sin\beta\\
			&\Leftrightarrow &  \sin\left[\left(\alpha+\beta\right)+\alpha\right]=3\sin\beta\\
			&\Leftrightarrow & \sin\left(\alpha+\beta\right)\cos\alpha+\sin\alpha\cos\left(\alpha+\beta\right)=3\sin\beta\\ 
			&\Leftrightarrow &  \dfrac{\sin\left(\alpha+\beta\right)}{\cos\left(\alpha+\beta\right)}\cos\alpha+\sin\alpha=\dfrac{3\sin\beta}{\cos\left(\alpha+\beta\right)}\\
			&\Leftrightarrow &  \dfrac{\sin\left(\alpha+\beta\right)}{\cos\left(\alpha+\beta\right)}=\dfrac{3\sin\beta}{\cos\alpha\cos\left(\alpha+\beta\right)}-\dfrac{\sin\alpha}{\cos\alpha}.\quad(*)
		\end{eqnarray*}
		Mà  $\dfrac{\sin\beta}{\cos\left(\alpha+\beta\right)}=\sin\alpha $, suy ra $(*)\Leftrightarrow\tan\left(\alpha+\beta\right)=\dfrac{3\sin\alpha}{\cos\alpha}-\dfrac{\sin\alpha}{\cos\alpha}=2\tan\alpha $.\\
		Vậy $\tan\left(\alpha+\beta\right)=2\tan\alpha $. 
	} 
\end{ex}

% nguồn: Ok_Toan11_TracNghiem_Lan2
\begin{ex}%[1D1V3-5]%[ma:0e8bb76b]
	Cho biết $\cos\alpha=-\dfrac{2}{3}$. Giá trị của  biểu thức $P=\dfrac{\cot\alpha+3\tan\alpha}{2\cot\alpha+\tan\alpha}$ bằng bao nhiêu?
	 \choice
		 {\True $P=\dfrac{19}{13}$}
		 {$P=\dfrac{25}{13}$}
		{$P=-\dfrac{25}{13}$} 
		{$P=-\dfrac{19}{13}$} 
	\loigiai{ 
		Ta có  $\cos\alpha=-\dfrac{2}{3}\Rightarrow{\tan ^2}\alpha=\dfrac{1}{\cos^2\alpha}-1=\dfrac{1}{\left(\dfrac{-2}{3}\right)^2}-1=\dfrac{5}{4}$.\\
		$P=\dfrac{\cot\alpha+3\tan\alpha}{2\cot\alpha+\tan\alpha}=\dfrac{\dfrac{1}{\tan\alpha}+3\tan\alpha}{\dfrac{2}{\tan\alpha}+\tan\alpha}=\dfrac{\dfrac{1+3\tan^2\alpha}{\tan\alpha}}{\dfrac{2+\tan^2\alpha}{\tan\alpha}}=\dfrac{1+3\tan^2\alpha}{2+\tan^2\alpha}=\dfrac{1+3\cdot\dfrac{5}{4}}{2+\dfrac{5}{4}}=\dfrac{19}{13}$. 
	 }
\end{ex}

% nguồn: Ok_Toan11_TracNghiem_Lan2
\begin{ex}%[1D1V3-5]%[ma:e93b9435]
	Nếu $5\sin\alpha=3\sin\left(\alpha+2\beta\right)$ thì 
	 \choice
		{$\tan\left(\alpha+\beta\right)=2\tan\beta  $}
		{$\tan\left(\alpha+\beta\right)=3\tan\beta $} 
		{\True $\tan\left(\alpha+\beta\right)=4\tan\beta  $}
		{$\tan\left(\alpha+\beta\right)=5\tan\beta  $}
	 \loigiai{
		 Ta có
		 \allowdisplaybreaks
		\begin{eqnarray*}
			&&5\sin\alpha=3\sin\left(\alpha+2\beta\right)\\ 
			&\Leftrightarrow&  5\sin\left[\left(\alpha+\beta\right)-\beta\right]=3\sin\left[\left(\alpha+\beta\right)+\beta\right]\\
			&\Leftrightarrow& 5\sin\left(\alpha+\beta\right)\cos\beta-5\cos\left(\alpha+\beta\right)\sin\beta=3\sin\left(\alpha+\beta\right)\cos\beta+3\cos\left(\alpha+\beta\right)\sin\beta  \\
			&\Leftrightarrow&  2\sin\left(\alpha+\beta\right)\cos\beta=8\cos\left(\alpha+\beta\right)\sin\beta \\ &\Leftrightarrow& \dfrac{\sin\left(\alpha+\beta\right)}{\cos\left(\alpha+\beta\right)}=4\dfrac{\sin\beta}{\cos\beta}\\ &\Leftrightarrow&\tan\left(\alpha+\beta\right)=4\tan\beta.
		\end{eqnarray*}
	 }
\end{ex}


%%%%%%%%%%%%%%%% ĐÚNG/SAI %%%%%%%%%%%%%%%%

% nguồn: Ok_dungsai_lop11_Lan2
\begin{ex}%[1D1H3-5]%[ma:44e22c29]
	Cho $\cos a=\dfrac{3}{4}$; $\sin a  > 0$; $\sin b=\dfrac{3}{5}$; $\cos b < 0$ . Xét tính đúng sai của các phát biểu sau
	 \choiceTF
		{\True Giá trị của $\tan  a=\dfrac{\sqrt 7}{3}$}
		{Giá trị của $\cot  b=-\dfrac{2}{3}$}
		{Giá  trị của $\cos 2a+\cos 2b$ thuộc khoảng $\left(\dfrac{1}{2};1\right)$}
		{Giá trị của $\cos\left(a+b\right)$ thuộc khoảng  $\left(-\dfrac{1}{2};-\dfrac{1}{3}\right)$}
	 \loigiai{
		Cho $\cos a=\dfrac{3}{4}$; $\sin a > 0$; $\sin b=\dfrac{3}{5}$; $\cos  b < 0$.
		\begin{itemchoice}
			\itemch  \textbf{Đúng.} \\
			Ta có $\tan ^2a+1=\dfrac{1}{\cos^2a}\Rightarrow\tan  a=\sqrt{\dfrac{1}{\cos^2a}-1}=\sqrt{\dfrac{16}{9}-1}=\dfrac{\sqrt 7}{3}$ .
			\itemch \textbf{Sai.}  \\
			Ta  có $\cot ^2b+1=\dfrac{1}{\sin^2b}\Rightarrow\cot b=-\sqrt{\dfrac{1}{\sin^2a}-1}=-\sqrt{\dfrac{25}{9}-1}=-\dfrac{4}{3}$.
			 \itemch \textbf{Sai.} \\
			Ta có $\cos 2a+\cos  2b=2\cos ^2a-1+1-2\sin ^2b=2.\dfrac{9}{16}-2.\dfrac{9}{25}=\dfrac{81}{200}=0,405\notin\left(\dfrac{1}{2};1\right)$.
			\itemch \textbf{Sai.}  \\
			Ta có $\heva{& \cos a=\dfrac{3}{4} \\ & \sin a>0 }\Rightarrow \sin a=\sqrt {1-{{\cos }^2}a}=\dfrac{\sqrt {7}}{4}$ và $\heva{& \sin b=\dfrac{3}{5} \\  & \cos b<0 }\Rightarrow \cos b=-\sqrt {1-{{\sin }^2}b}=-\dfrac{4}{5}$.\\
			Suy ra $\cos \left( a+b \right)=\cos a\cos b-\sin a\sin b=\frac{3}{4}.\left( -\frac{4}{5} \right)-\frac{\sqrt {7}}{4}.\frac{3}{5}=-\frac{3}{5}\left( 1+\frac{\sqrt {7}}{4} \right)\notin ( -\dfrac{1}{2};-\dfrac{1}{3})$. 
		\end{itemchoice}
	 }
\end{ex}

% nguồn: Ok_dungsai_lop11_Lan2
\begin{ex}%[1D1H3-5]%[ma:4abf121a]
	 Biết $\sin 2\alpha=-\dfrac{4}{5},\dfrac{\pi}{2}<\alpha <\dfrac{3\pi}{2}$. Khi đó
	 \choiceTF
		 {\True $\cos\alpha < 0$}
		{\True $  2\sin\alpha\cos\alpha=-\dfrac{4}{5}$}
		{\True $\cos\alpha=\dfrac{-2}{\sqrt 5},\sin\alpha=\dfrac{1}{\sqrt 5}$} 
		 {$\cos\alpha=\dfrac{-1}{\sqrt 5},\sin\alpha=-\dfrac{2}{\sqrt 5}$}
	\loigiai{ 
		Vì $\dfrac{\pi}{2}<\alpha  <\dfrac{3\pi}{2}$ nên $\cos\alpha < 0$. Ta có hệ phương trình\\
		${\heva{&\sin ^2 \alpha+\cos ^2 \alpha=1 \\ &2 \sin \alpha \cos \alpha=-\dfrac{4}{5}}}$ 
		$\Rightarrow \left\{{\begin{array}{*{20}{l}}
		{\dfrac{4}{25  c o s{^2}\alpha}+c o s{^2}\alpha=1}\\
		{s i n\alpha=-\dfrac{2}{5 c  o s\alpha}}
	\end{array} \Rightarrow \left\{{\begin{array}{*{20}{l}}
	 {25\cos^4\alpha-25\cos^2\alpha+4=0}\\
	 {\sin\alpha=-\dfrac{2}{5\cos\alpha}}
	\end{array}}\right.}\right.$\\
	$\Rightarrow\left\{\begin{array}{l}
	\left[\begin{array}{l}
	\cos{^2}\alpha=\dfrac{4}{5}\\ 
	 \cos{^2}\alpha=\dfrac{1}{5}
	\end{array}\right.\\
	\sin\alpha=-\dfrac{2}{5\cos\alpha} 
	\end{array}\right.\Rightarrow\left\{\begin{array}{l}
	\left[\begin{array}{l}
	\cos\alpha=\dfrac{-2}{\sqrt  5}\\
	\cos\alpha=\dfrac{-1}{\sqrt 5} 
	\end{array}\right.\\
	\sin\alpha=-\dfrac{2}{5\cos\alpha} 
	\end{array}\right. \Rightarrow \hoac{& \cos \alpha =\dfrac{-2}{\sqrt {5}},\sin \alpha =\dfrac{1}{\sqrt {5}} \\ & \cos \alpha =\dfrac{-1}{\sqrt {5}},\sin \alpha =\dfrac{2}{\sqrt {5}}}$}
\end{ex}

% nguồn: Ok_dungsai_lop11_Lan2
\begin{ex}%[1D1H3-5]%[ma:bb440ab8]
	Cho biết $\tan x=\sqrt 2 $ và $ 0 < x <90^\circ$. Khi  đó
	 \choiceTF
		{\True $\cos x > 0$} 
		{\True $\cos  x=\dfrac{\sqrt 3}{3}$}
		{\True $\sin  x=\dfrac{\sqrt 6}{3}$}
		{$\cos\left(x-30^\circ\right)=\dfrac{3-\sqrt 6}{6}$} 
	\loigiai{ 
		Vì $ 0 < x <90^\circ$ nên $\cos x >  0$.\\
		Ta có $\dfrac{1}{\cos^2x}=1+\tan ^2x=1+2=3\Rightarrow{\cos ^2}x=\dfrac{1}{3}\Rightarrow\cos x=\dfrac{\sqrt 3}{3}$.\\ 
		Măt khác $\tan  x=\dfrac{\sin x}{\cos x}\Rightarrow\sin x=\tan x\cos x=\dfrac{\sqrt 6}{3}$.\\
	Khi đó $\cos\left(x-30^\circ\right)=\cos x\cos{30^\circ}+\sin  x\sin{30^\circ}=\dfrac{\sqrt 3}{3}\cdot\dfrac{\sqrt 3}{2}+\dfrac{\sqrt 6}{3}\cdot\dfrac{1}{2}=\dfrac{3+\sqrt 6}{6}$.}
\end{ex}

% nguồn: Ok_dungsai_lop11_Lan2
\begin{ex}%[1D1H3-5]%[ma:5bd8ceab]
	Xét tính đúng-sai của các phát  biểu sau
	 \choiceTF
		{\True $\sin\left(a+b\right)=\sin a\cos b+\cos  a\sin b$}
		{$\sin{31^0}.\cos{12^0}+\cos{12^0}.\sin{31^0}=\sin{19^0}$} 
		{\True Cho $\cos x=\dfrac{4}{5}, x\in ( -\dfrac{\pi }{2};0)$. Giá trị của $\sin 2x$  là $-\dfrac{24}{25}$}
		 {Cho $\sin \alpha =\dfrac{1}{3}$ và $\dfrac{\pi }{2}<\alpha <\pi $. Biết giá trị của $\cos ( \alpha -\dfrac{\pi }{6})=\dfrac{1-a\sqrt {6}}{b}$ với $ a,b\in \mathbb{N}$ thì $ a+b=4$}
	\loigiai{ 
	\begin{itemchoice}
		\itemch \textbf{Đúng.}  \\
		Theo  công thức cộng, ta có $\sin\left(a+b\right)=\sin a\cos b+\cos a\sin b$.
		 \itemch \textbf{Sai.} \\
		Áp dụng công thức cộng, ta  có
		$\sin{31^0}.\cos{12^0}+\cos{12^0}.\sin{31^0}=\sin\left(31^0+12^0\right)=\sin{43^0}$. 
		 \itemch \textbf{Đúng.} \\
		Ta có $\sin ^2x=1-\cos ^2x=1-\dfrac{16}{25}=\dfrac{9}{25}$ $\Rightarrow\sin x=-\dfrac{3}{5}$. Mặt khác, vì $x\in\left(-\dfrac{\pi}{2};0\right)\Rightarrow\sin x < 0$.\\ 
		 Vậy $\sin 2x=2\sin x.\cos x=2.\dfrac{4}{5}.\left(-\dfrac{3}{5}\right)=-\dfrac{24}{25}$.
		 \itemch \textbf{Sai.} \\
		Vì $\sin\alpha=\dfrac{1}{3},\dfrac{\pi}{2}<\alpha <\pi $ nên  $\cos\alpha=-\dfrac{2\sqrt 2}{3}$.\\
		Do đó $\cos\left(\alpha-\dfrac{\pi}{6}\right)=\cos\alpha .\cos\dfrac{\pi}{6}+\sin\alpha .\sin\dfrac{\pi}{6}=-\dfrac{2\sqrt  2}{3}.\dfrac{\sqrt 3}{2}+\dfrac{1}{3}.\dfrac{1}{2}=\dfrac{1-2\sqrt 6}{6}$.\\
		 Khi đó $ a=2,\,\,b=6$, ta có $ a+b=8$.
	\end{itemchoice}
	 }
\end{ex}

% nguồn: Ok_dungsai_lop11_Lan2
\begin{ex}%[1D1H3-5]%[ma:7cece702]
	Xét tính  đúng-sai của các đẳng thức sau
	\choiceTF 
		{$\sin\left(x-\dfrac{\pi}{6}\right)=\dfrac{\sqrt  3}{2}\sin x+\dfrac{1}{2}\cos x$}
		 {\True $4\sin x\cdot\sin 2x\cdot\sin 3x=\sin 4x+\sin 2x-\sin 6x$}
		{$ 1+\sin 2x+\cos 2x=2\sqrt  2\sin x\cos\left(x-\dfrac{\pi}{4}\right)$}
		{$\sin ^2x+2\sin\left(a-x\right)\cdot\sin x\cdot\cos a+\sin  ^2\left(a-x\right)=\cos ^2a$}
	\loigiai{ 
		\begin{itemchoice}
			\itemch  \textbf{Sai.} \\
			Theo công thức cộng ta có $\sin\left(x-\dfrac{\pi}{6}\right)=\sin x\cdot\cos\dfrac{\pi}{6}-\cos x\cdot\sin\dfrac{\pi}{6}=\dfrac{\sqrt  3}{2}\sin x-\dfrac{1}{2}\cos x$.\\
			Do đó  $\sin\left(x-\dfrac{\pi}{6}\right)=\dfrac{\sqrt 3}{2}\sin x+\dfrac{1}{2}\cos x$ là sai.
			\itemch \textbf{Đúng.}  \\
			Áp dụng quy  tắc từ tích sang tổng, ta được\\
			$4\sin x\cdot\sin 2x\cdot\sin 3x=2\left(\cos x-\cos 3x\right)\sin 3x=\sin 4x+\sin  2x-\sin 6x$.\\
			Do đó $4\sin x\cdot\sin 2x\cdot\sin 3x=\sin 4x+\sin  2x-\sin 6x$ là đúng.
			 \itemch \textbf{Sai.} \\
			Ta có $ 1+\sin 2x+\cos 2x=2\sin x\cos x+2\cos ^2x=2\cos x(\sin x+\cos x)$.\\ 
			Mà  $\sin x+\cos x=\sqrt 2\left(\cos x\dfrac{1}{\sqrt 2}+\sin x\dfrac{1}{\sqrt 2}\right)=\sqrt 2\cos\left(x-\dfrac{\pi}{4}\right).$\\
			Vậy $ 1+\sin 2x+\cos  2x=2\sqrt 2\sin x\cdot\cos\left(x-\dfrac{\pi}{4}\right)$ là sai.
			\itemch \textbf{Sai.}  \\
			Ta  có
			$\sin ^2x+2\sin\left(a-x\right)\cdot\sin x\cdot\cos a+\sin  ^2\left(a-x\right)$\\
			$=\sin ^2x+\sin\left(a-x\right)\cdot\left[\sin\left(x+a\right)+\sin\left(x-a\right)\right]+\sin  ^2\left(a-x\right)$\\
			 $=\sin ^2x+\sin\left(a-x\right)\cdot\left[\sin\left(x+a\right)+\sin\left(x-a\right)+\sin\left(a-x\right)\right]$\\
			$=\sin ^2x+\sin\left(a-x\right)\cdot\left[\sin\left(x+a\right)+\sin\left(x-a\right)-\sin\left(x-a\right)\right]$\\ 
			$=\sin ^2x+\sin\left(a-x\right)\cdot\sin\left(x+a\right)$\\ 
			$=\dfrac{1-\cos 2x}{2}-\dfrac{\cos 2a-\cos 2x}{2}=\dfrac{1-\cos 2a}{2}\,=\sin ^2a.$\\ 
			Do đó $\sin ^2x+2\sin\left(a-x\right)\cdot\sin x\cos a+\sin ^2\left(a-x\right)=\cos  ^2a$ là sai.
		\end{itemchoice}
	 }
\end{ex}

% nguồn: Ok_dungsai_lop11_Lan2
\begin{ex}%[1D1H3-5]%[ma:bb9ba528]
	Cho biết  $\sin\alpha=-\dfrac{12}{13},\dfrac{3\pi}{2}<\alpha < 2\pi $. Khi đó
	 \choiceTF
		{\True $\cos\alpha > 0$} 
		{\True $\cos\alpha=\dfrac{5}{13}$} 
		{\True $\tan\alpha=-\dfrac{12}{5}$} 
		{$\cos\left(\dfrac{\pi}{3}-\alpha\right)=\dfrac{5-\sqrt 3}{26}$} 
	 \loigiai{
		Ta có $\cos\left(\dfrac{\pi}{3}-\alpha\right)=\cos\dfrac{\pi}{3}\cos\alpha+\sin\dfrac{\pi}{3}\sin\alpha = \dfrac{1}{2}\cos\alpha+\dfrac{\sqrt 3}{2}\cdot\left(-\dfrac{12}{13}\right)=\dfrac{1}{2}\cos\alpha-\dfrac{6\sqrt  3}{13}$.\\
		Vì $\dfrac{3\pi}{2}<\alpha  < 2\pi\Rightarrow\cos\alpha > 0$.\\
		Mặt khác $\cos\alpha=\sqrt{1-\sin^2\alpha}=\dfrac{5}{13}$. Suy ra $\cos\left(\dfrac{\pi}{3}-\alpha\right)=\dfrac{1}{2}\cdot\dfrac{5}{13}-\dfrac{6\sqrt 3}{13}=\dfrac{5-12\sqrt 3}{26}$. 
	} 
\end{ex}

% nguồn: Ok_dungsai_lop11_Lan2
\begin{ex}%[1D1H3-5]%[ma:b098d45d]
	Cho $\cot x=-\sqrt 3, \dfrac{3\pi}{2}< x < 2\pi  $. Khi đó
	\choiceTF 
		{\True  $\sin x=-\dfrac{\sqrt{10}}{10}$}
		{$\cos  x=\dfrac{\sqrt 3}{10}$}
		 {\True $\sin\left(\dfrac{4\pi}{3}-x\right)=\dfrac{-\sqrt{10}}{5}$}
		{\True $\tan\left(x+\dfrac{\pi}{3}\right)=\dfrac{\sqrt  3}{3}$}
	\loigiai{ 
		Ta  có
		\begin{itemize}
			\item $\tan x=\dfrac{1}{\cot x}=-\dfrac{\sqrt  3}{3}$.
			\item $\dfrac{1}{\sin^2x}=1+\cot  ^2x=1+(-\sqrt 3)^2=10\Rightarrow{\sin ^2}x=\dfrac{1}{10}\Rightarrow \sin x=\pm\dfrac{\sqrt{10}}{10}$.\\
			 Vì $\dfrac{3\pi}{2}< x < 2\pi$ nên $\sin x=-\dfrac{\sqrt{10}}{10}$.
			\item $\cos  x=\cot x\cdot\sin x=-\sqrt 3 .\left(-\dfrac{\sqrt{10}}{10}\right)=\dfrac{\sqrt{30}}{10}$.
			\item $\sin\left(\dfrac{4\pi}{3}-x\right)=\sin\dfrac{4\pi}{3}\cos x-\cos\dfrac{4\pi}{3}\sin x=-\dfrac{\sqrt 3}{2}\cdot\left(\dfrac{\sqrt{30}}{10}\right)-\dfrac{-1}{2}\cdot\dfrac{-\sqrt{10}}{10}=\dfrac{-\sqrt{10}}{5}$. 
			 \item $\tan\left(x+\dfrac{\pi}{3}\right)=\dfrac{\tan x+\tan\dfrac{\pi}{3}}{1-\tan x\cdot\tan\dfrac{\pi}{3}}=\dfrac{\dfrac{-\sqrt 3}{3}+\sqrt 3}{1-\dfrac{-\sqrt 3}{3}\cdot\sqrt 3}=\dfrac{\sqrt 3}{3}$.
		\end{itemize}
	} 
\end{ex}

% nguồn: Ok_dungsai_lop11_Lan2
\begin{ex}%[1D1H3-5]%[ma:6319dd8a]
	Cho biết $\sin\alpha=\dfrac{1}{3}$ và $\dfrac{\pi}{2}<\alpha <\pi $. Khi  đó
	 \choiceTF
		 {\True $\cos\alpha=-\dfrac{2\sqrt 2}{3}$}
		{\True $\sin  2\alpha=-\dfrac{4\sqrt 2}{9}$}
		{\True  $\cos 2\alpha=\dfrac{7}{9}$}
		{$\cot 2\alpha=\dfrac{7\sqrt 2}{8}$} 
	\loigiai{ 
		Vì $\dfrac{\pi}{2}<\alpha<\pi$ nên $\cos \alpha<0$. Ta có\\ 
		\begin{itemize}
			\item  $\cos^2\alpha=1-\sin^2\alpha=1-\dfrac{1}{9}=\dfrac{8}{9}\Rightarrow\cos\alpha=-\dfrac{2\sqrt 2}{3}$.
			\item $\sin2\alpha=2\sin\alpha\cos\alpha=2\cdot\dfrac{1}{3}\cdot\left(-\dfrac{2\sqrt  2}{3}\right)=-\dfrac{4\sqrt 2}{9}$.
			\item $\cos 2\alpha=1-2\sin ^2\alpha=1-2\left(\dfrac{1}{3}\right)^2=\dfrac{7}{9}$. 
			\item $\tan 2\alpha=\dfrac{\sin 2\alpha}{\cos 2\alpha}=-\dfrac{4\sqrt 2}{7}; \cot 2\alpha=\dfrac{1}{\tan 2\alpha}=-\dfrac{7\sqrt  2}{8}$.
		\end{itemize}
	} 
\end{ex}

% nguồn: Ok_dungsai_lop11_Lan2
\begin{ex}%[1D1H3-5]%[ma:b75b6ed6]
	Cho  biết $\cos 2\alpha=-\dfrac{1}{4}$ và $\pi <\alpha <\dfrac{3\pi}{2}$. Khi đó
	 \choiceTF
		{\True  $\sin\alpha < 0,\cos\alpha < 0$}
		 {$\sin\alpha=\dfrac{\sqrt{10}}{4}$}
		 {$\cos\alpha=\dfrac{\sqrt 6}{4}$}
		{\True  $\cot\alpha=\dfrac{\sqrt{15}}{5}$}
	 \loigiai{
		Vì $\pi <\alpha <\dfrac{3\pi}{2}$ nên  $\sin\alpha < 0, \cos\alpha < 0$. Ta có\\
		\begin{itemize}
			\item $\cos 2\alpha=-\dfrac{1}{4}\Rightarrow  1-2\sin ^2\alpha=-\dfrac{1}{4}\Rightarrow{\sin^2}\alpha=\dfrac{5}{8}\Rightarrow\sin\alpha=-\dfrac{\sqrt{10}}{4}$.
			\item $\cos\alpha=-\sqrt{1-\sin^2\alpha}=-\sqrt{1-\dfrac{10}{16}}=-\dfrac{\sqrt 6}{4}$. 
			\item  $\tan\alpha=\dfrac{\sin\alpha}{\cos\alpha}=\dfrac{\sqrt{15}}{3}$.
			 \item $\cot\alpha=\dfrac{1}{\tan\alpha}=\dfrac{\sqrt{15}}{5}$
		\end{itemize}
	 }
\end{ex}

% nguồn: Ok_dungsai_lop11_Lan2
\begin{ex}%[1D1H3-5]%[ma:7f347f32]
	 Cho biết $\cos x=-\dfrac{12}{13}$ và $\pi < x <\dfrac{3\pi}{2}$, khi đó
	\choiceTF 
		{$\sin x  > 0$}
		 {\True $\sin x=-\dfrac{5}{13}$}
		 {$\cot x=\dfrac{5}{12}$}
		 {\True $\sin\left(\dfrac{\pi}{3}-x\right)=\dfrac{5-12\sqrt 3}{26}$}
	\loigiai{ 
		Vì $\pi < x  <\dfrac{3\pi}{2}$ nên $\sin x < 0$. Ta có\\
		\begin{itemize}
			\item  $\cos x=-\dfrac{12}{13}$ $\Rightarrow $ $\sin x=-\sqrt{1-\cos^2x}=-\sqrt{1-\left(-\dfrac{12}{13}\right)^2}=-\dfrac{5}{13}$ $\Rightarrow $ $\cot x=\dfrac{\cos x}{\sin x}=\dfrac{12}{5}$.
			\item $\sin\left(\dfrac{\pi}{3}-x\right)=\sin\dfrac{\pi}{3}\cos x-\cos\dfrac{\pi}{3}\sin x=\dfrac{\sqrt 3}{2}\left(-\dfrac{12}{13}\right)-\dfrac{1}{2}\left(-\dfrac{5}{13}\right)=\dfrac{5-12\sqrt  3}{26}$.
		\end{itemize}
	} 
\end{ex}

% nguồn: Ok_dungsai_lop11_Lan2
\begin{ex}%[1D1H3-5]%[ma:55c27e77]
	Biết $\sin a=\dfrac{8}{17},\tan b=\dfrac{5}{12}$ và $ a$, $ b$ là các góc  nhọn. Khi đó
	 \choiceTF
		{\True $\tan  a=\dfrac{8}{15}$}
		 {\True $\sin (a-b)=\dfrac{21}{221}$}
		 {$\cos (a+b)=\dfrac{14}{22}$}
		 {$\tan (a+b)=\dfrac{17}{14}$}
	\loigiai{ 
		Vì $ a,b$ là các  góc nhọn nên $\cos a > 0,\cos b > 0$.\\
		Ta  có
		\begin{itemize}
			\item $\cos a=\sqrt{1-\sin^2a}=\dfrac{15}{17}  \Rightarrow \tan a=\dfrac{\sin a}{\cos a}=\dfrac{8}{15}$.
			\item $\cos b=\sqrt{\dfrac{1}{1+\tan^2b}}=\dfrac{12}{13}\Rightarrow\sin  b=\cos b\tan b=\dfrac{5}{13}$.
		\end{itemize}
		 Khi đó
		\begin{itemize}
			\item $\sin  (a-b)=\sin a\cos b-\cos a\sin b=\dfrac{8}{17}\cdot\dfrac{12}{13}-\dfrac{15}{17}\cdot\dfrac{5}{13}=\dfrac{21}{221}$.
			\item $\cos (a+b)=\cos a\cos b-\sin a\sin  b=\dfrac{15}{17}\cdot\dfrac{12}{13}-\dfrac{8}{17}\cdot\dfrac{5}{13}=\dfrac{140}{221}$.
			\item $\tan (a+b)=\dfrac{\tan a+\tan b}{1-\tan  a\tan b}=\dfrac{\dfrac{8}{15}+\dfrac{5}{12}}{1-\dfrac{8}{15}\cdot\dfrac{5}{12}}=\dfrac{171}{140}$.
		\end{itemize}
	 }
\end{ex}

% nguồn: Ok_dungsai_lop11_Lan2
\begin{ex}%[1D1H3-5]%[ma:2ff7c3b8]
	Cho $\cos\left(\dfrac{\pi}{2}+x\right)=-\dfrac{1}{3}$, với $x\in\left(\dfrac{\pi}{2}\,;\,\pi\right)$. Xét tính  đúng sai của các mệnh đề sau
	 \choiceTF
		{\True $\sin 2x=-\dfrac{4\sqrt  2}{9}$}
		{$\cos 2x=-\dfrac{4\sqrt  2}{9}$}
		{\True  $\tan 2x=-\dfrac{4\sqrt 2}{7}$}
		 {$\sin x+\cos x=1$}
	\loigiai{ 
		Cho  ${cos}\left(\dfrac{\pi}{2}+x\right)=-\dfrac{1}{3} \Rightarrow \sin x=\dfrac{1}{3}$.\\
		Với $x\in\left(\dfrac{\pi}{2}\,;\,\pi\right)$ ta có  $\cos x < 0\Rightarrow\cos x=-\sqrt{1-\sin^2x}=-\sqrt{1-\left(\dfrac{1}{3}\right)^2}=-\dfrac{2\sqrt 2}{3}$.\\
		\begin{itemchoice}
			\itemch \textbf{Đúng.} \\ 
			Vì $\sin 2x=2\cdot\sin x\cdot\cos x=2\cdot\dfrac{1}{3}\cdot\left(-\dfrac{2\sqrt  2}{3}\right)=-\dfrac{4\sqrt 2}{9}$.
			\itemch \textbf{Sai.}  \\
			$\cos  2x=1-2\sin ^2x=1-2\left(\dfrac{1}{3}\right)^2=\dfrac{7}{9}$.
			\itemch  \textbf{Đúng.} \\
			$\tan 2x=\dfrac{\sin 2x}{\cos 2x}=\dfrac{-4\sqrt  2}{7}$.
			\itemch \textbf{Sai.} \\ 
			$\sin  x+\cos x=\dfrac{1}{3}-\dfrac{2\sqrt 2}{3}=\dfrac{1-2\sqrt 2}{3}.$
		\end{itemchoice}
	 }
\end{ex}

% nguồn: Ok_dungsai_lop11_Lan2
\begin{ex}%[1D1H3-5]%[ma:81be90af]
	Cho biết $\sin\alpha=\dfrac{3}{5},\dfrac{\pi}{2}<\alpha <\pi $. Khi  đó
	 \choiceTF
		{\True  $\cos\alpha < 0$}
		{\True $\cos\alpha=-\dfrac{4}{5}$} 
		{$\tan\alpha=\dfrac{3}{4}$} 
		{$\tan\left(\alpha+\dfrac{\pi}{3}\right)=\dfrac{48-\sqrt 3}{11}$} 
	\loigiai{ 
		Vì $\dfrac{\pi}{2}<\alpha <\pi\Rightarrow\cos\alpha  < 0$.\\
		Ta  có $\cos\alpha=-\sqrt{1-\sin^2\alpha}=-\dfrac{4}{5}.\\
		\Rightarrow \tan\alpha=\dfrac{\sin\alpha}{\cos\alpha}=-\dfrac{3}{4}$,  $
		\tan\left(\alpha+\dfrac{\pi}{3}\right)=\dfrac{\tan\alpha+\tan\dfrac{\pi}{3}}{1-\tan\alpha\tan\dfrac{\pi}{3}}=\dfrac{\tan\alpha+\sqrt  3}{1-\sqrt 3\tan\alpha}$.\\
		$\Rightarrow \tan\left(\alpha+\dfrac{\pi}{3}\right)=\dfrac{-\dfrac{3}{4}+\sqrt  3}{1-\sqrt 3\left(-\dfrac{3}{4}\right)}=\dfrac{-3+4\sqrt 3}{4+3\sqrt 3}=\dfrac{48-25\sqrt 3}{11}$.
	 }
\end{ex}

% nguồn: Ok_dungsai_lop11_Lan2
\begin{ex}%[1D1H3-5]%[ma:9fe42cb6]
	Biết  $\sin\alpha=\dfrac{1}{3}$ và $ 0 <\alpha <\dfrac{\pi}{2}$. Khi đó
	 \choiceTF
		{$\sin 2\alpha=\dfrac{\sqrt  2}{9}$}
		{\True $\cos  2\alpha=\dfrac{7}{9}$}
		{$\tan 2\alpha=\dfrac{\sqrt  2}{7}$}
		{$\cot 2\alpha=\dfrac{7\sqrt 2}{2}$} 
	\loigiai{ 
		Ta có $\cos  ^2\alpha+\sin ^2\alpha =1 \Rightarrow \cos ^2\alpha=1-\sin ^2\alpha=1-\left(\dfrac{1}{3}\right)^2=\dfrac{8}{9}\Rightarrow\cos\alpha=\pm\dfrac{2\sqrt 2}{3}$.\\
		Vì $ 0 <\alpha  <\dfrac{\pi}{2}$ nên $\cos\alpha=\dfrac{2\sqrt 2}{3}$. Khi đó
		\begin{itemize}
			\item  $\sin 2\alpha=2\sin\alpha\cdot\cos\alpha=2\cdot\dfrac{1}{3}\cdot\dfrac{2\sqrt 2}{3}=\dfrac{4\sqrt 2}{9}$.
			\item  $\cos 2\alpha=\cos ^2\alpha-\sin ^2\alpha=\left(\dfrac{2\sqrt 2}{3}\right)^2-\left(\dfrac{1}{3}\right)^2=\dfrac{7}{9}$.
			\item $\tan 2\alpha=\dfrac{\sin 2\alpha}{\cos 2\alpha}=\dfrac{4\sqrt 2}{2}:\dfrac{7}{9}=\dfrac{4\sqrt  2}{7}$.
			\item $\cot 2\alpha=\dfrac{\cos 2\alpha}{\sin 2\alpha}=\dfrac{4\sqrt  2}{2}:\dfrac{7}{9}=\dfrac{7\sqrt 2}{8}$.
		\end{itemize}
	 }
\end{ex}

% nguồn: Ok_dungsai_lop11_Lan2
\begin{ex}%[1D1H3-5]%[ma:f4f1f2e2]
	 Cho $\sin\alpha=\dfrac{2}{3}, \dfrac{\pi}{2}<\alpha <\pi $. Khi đó
	\choiceTF 
		 {\True $\cos\alpha=-\dfrac{\sqrt 5}{3}$}
		{\True $\tan\alpha=-\dfrac{2\sqrt  5}{5}$}
		 {$\cos\left(\dfrac{\pi}{3}+\alpha\right)=\dfrac{\sqrt 5-2\sqrt 3}{6}$}
		 {$\cos\left(\dfrac{\pi}{4}-\alpha\right)=\dfrac{\sqrt{10}-2\sqrt 2}{6}$}
	 \loigiai{
		Ta có $\cos ^2\alpha+\sin ^2\alpha  =1$ $\Rightarrow \cos ^2\alpha=1-\sin ^2\alpha=1-\left(\dfrac{2}{3}\right)^2=\dfrac{5}{9}\Rightarrow\cos\alpha=\pm\dfrac{\sqrt 5}{3}$.\\
		Vì $\dfrac{\pi}{2}<\alpha <\pi $  nên $\cos\alpha=-\dfrac{\sqrt 5}{3}$. Khi đó
		\begin{itemize}
			 \item $\tan\alpha=\dfrac{\sin\alpha}{\cos\alpha}=-\dfrac{2\sqrt 5}{5}$.
			\item $\cos\left(\dfrac{\pi}{3}+\alpha\right)=\cos\dfrac{\pi}{3}\cos\alpha-\sin\dfrac{\pi}{3}\sin\alpha=\dfrac{1}{2}\cdot\left(\dfrac{-\sqrt 5}{3}\right)-\dfrac{\sqrt 3}{2}\cdot\dfrac{2}{3}=\dfrac{-\sqrt 5-2\sqrt 3}{6}$. 
			\item  $\cos\left(\dfrac{\pi}{4}-\alpha\right)=\cos\dfrac{\pi}{4}\cos\alpha+\sin\dfrac{\pi}{4}\sin\alpha=\dfrac{\sqrt 2}{2}\cdot\left(\dfrac{-\sqrt 5}{3}\right)+\dfrac{\sqrt 2}{2}\cdot\dfrac{2}{3}=\dfrac{-\sqrt{10}+2\sqrt 2}{6}$.
		\end{itemize}
	} 
\end{ex}

% nguồn: Ok_dungsai_lop11_Lan2
\begin{ex}%[1D1H3-5]%[ma:721cf5f0]
	Biết $ 0 < a,b <\dfrac{\pi}{2},a+b=\dfrac{\pi}{4}$ và  $\tan a\tan b=3-2\sqrt 2 $. Khi đó
	 \choiceTF
		{\True  $\tan a+\tan b=-2+2\sqrt 2\rm$}
		 {\True $\tan a=-1+\sqrt 2 $}
		{$\tan b=-1-\sqrt 2 $} 
		{$\tan a-\tan b=-2-2\sqrt  2\rm$}
	 \loigiai{
		Ta có  $\tan (a+b)=\dfrac{\tan a+\tan b}{1-\tan a\tan b}\Rightarrow\tan a+\tan b=\tan (a+b)(1-\tan a\tan b)$.\\
		 Mà $ a+b=\dfrac{\pi}{4}$ và $\tan a\tan b=3-2\sqrt 2$ $\Rightarrow\tan a+\tan b=\tan\dfrac{\pi}{4}[1-(3-2\sqrt 2)]=-2+2\sqrt 2$.\\
		Đặt $ S=\tan a+\tan b, P=\tan a\tan  b$.\\
		Khi  đó $\tan a,\tan b$ là nghiệm của phương trình $X^2-SX+P=0\\
		\Leftrightarrow{X^2}-(2\sqrt 2-2)X+3-2\sqrt 2=0$\\ 
		$\Leftrightarrow  X=-1+\sqrt 2\rm$. Suy ra $\tan a=\tan b=-1+\sqrt 2 $.
	} 
\end{ex}

% nguồn: Ok_dungsai_lop11_Lan2
\begin{ex}%[1D1H3-5]%[ma:4f7b140b]
	 Xác định tính đúng sai của các mệnh đề dưới đây
	\choiceTF 
		{\True $\sin 2x=2\sin x\cos  x$}
		{$\cos a-\cos  b=2\sin\dfrac{a+b}{2}\sin\dfrac{a-b}{2}$}
		{$\tan\left(a+b\right)+\tan  a+\tan b=\tan\left(a+b\right)\cdot\tan a\cdot\tan b$}
		{\True Tam giác $ ABC$ cân tại  $ C$ khi và chỉ khi $\dfrac{\sin A+\sin B}{\cos A+\cos B}=\dfrac{1}{2}\left(\tan A+\tan B\right)$}
	 \loigiai{
		\begin{itemchoice}
			\itemch  \textbf{Đúng.} \\
			Mệnh  đề a) đúng theo lý thuyết.
			\itemch  \textbf{Sai.} \\
			Mệnh đề b) sai  do $\cos a-\cos b=-2\sin\dfrac{a+b}{2}\sin\dfrac{a-b}{2}$.
			\itemch \textbf{Sai.}  \\
			 Mệnh đề c) sai vì\\
			 $\tan\left(a+b\right)+\tan a+\tan b=\tan\left(a+b\right)+\tan\left(a+b\right)\left(1-\tan a.\tan b\right)$\\
			$=2\tan\left(a+b\right)-\tan\left(a+b\right)\cdot\tan a\cdot\tan  b\ne\tan\left(a+b\right)\cdot\tan a\cdot\tan b$.
			\itemch \textbf{Đúng.} \\ 
			Mệnh  đề d) đúng vì\\
			\textbf{Điều kiện cần:}  Giả sử tam giác $ABC$ thỏa mãn điều kiện $\dfrac{\sin A+\sin B}{\cos A+\cos B}=\dfrac{1}{2}\left(\tan A+\tan B\right)\,\,(3)$ \\
			Ta chứng minh tam giác $ABC$ cân. Thật vậy, ta có\\ 
			 $(3)\Rightarrow\dfrac{\sin A+\sin B}{\cos A+\cos B}=\dfrac{1}{2}\dfrac{\sin A\cos B+\sin B\cos A}{\cos A\cos B}$ $\\
			\Rightarrow\dfrac{2\cdot\sin\dfrac{A+B}{2}\cdot\cos\dfrac{A-B}{2}}{2\cdot\cos\dfrac{A+B}{2}\cdot\cos\dfrac{A-B}{2}}=\dfrac{1}{2}\dfrac{\sin A\cos B+\sin B\cos A}{\cos A\cos  B}$\\
			$\Rightarrow\dfrac{\sin\left(\dfrac{\pi}{2}-\dfrac{C}{2}\right)}{\cos\left(\dfrac{\pi}{2}-\dfrac{C}{2}\right)}=\dfrac{1}{2}\cdot\dfrac{\sin\left(A+B\right)}{\cos A\cdot\cos B}$\\ 
			$\Rightarrow\dfrac{\cos\dfrac{C}{2}}{\sin\dfrac{C}{2}}=\dfrac{1}{2}\dfrac{\sin C}{\cos A\cdot\cos B}$$\Rightarrow\dfrac{\cos\dfrac{C}{2}}{\sin\dfrac{C}{2}}=\dfrac{2\sin\dfrac{C}{2}\cos\dfrac{C}{2}}{2\cos A\cdot\cos B}$$\Rightarrow\cos  A\cdot\cos B=\sin ^2\dfrac{C}{2}$\\
			$\Rightarrow  2\cos A\cdot\cos B=1-\cos C$\\
			 $\Rightarrow 2\cos A\cdot\cos B=1+\cos\left(A+B\right)$\\
			$\Rightarrow 2\cos A\cdot\cos B=1+\cos  A\cdot\cos B-\sin A\cdot\sin B$\\
			 $\Rightarrow\cos\left(A-B\right)=1$\\
			$\Rightarrow  A=B$\\
			$\Rightarrow $Tam  giác $ABC\,$ cân tại $C$.\\
			\textbf{Điều kiện đủ} Giả sử tam giác $ABC$ cân tại đỉnh  $C$ $\Rightarrow A=B\Rightarrow(3)$ đúng.
		\end{itemchoice}
	} 
\end{ex}

% nguồn: Ok_dungsai_lop11_Lan2
\begin{ex}%[1D1H3-5]%[ma:1db05742]
	Biết $\sin  a=\dfrac{8}{17}$, $\tan b=\dfrac{5}{12}$ và $ a$, $ b$ là các góc nhọn. Khi đó
	 \choiceTF
		{\True $\tan  a=\dfrac{8}{15}$}
		 {\True $\sin (a-b)=\dfrac{21}{221}$}
		{$\cos (a+b)=\dfrac{14}{22}$} 
		{$\tan  2(a+b)=\dfrac{171}{140}$}
	 \loigiai{
		\begin{itemchoice}
			 \itemch \textbf{Đúng.} \\
			Vì $ a,b$ là các góc nhọn nên $\cos a > 0,\cos b >  0$.\\
			Ta có $\cos a=\sqrt{1-\sin^2a}=\dfrac{15}{17}\Rightarrow\tan a=\dfrac{\sin  a}{\cos a}=\dfrac{8}{15}$ nên là mệnh đề đúng.
			\itemch  \textbf{Đúng.} \\
			$\cos b=\sqrt{\dfrac{1}{1+\tan^2b}}=\dfrac{12}{13}\Rightarrow\sin  b=\cos b\tan b=\dfrac{5}{13}{\rm}$.\\
			Khi đó $\sin (a-b)=\sin a\cos b-\cos a\sin b=\dfrac{8}{17}\cdot\dfrac{12}{13}-\dfrac{15}{17}\cdot\dfrac{5}{13}=\dfrac{21}{221}$ nên  là mệnh đề đúng.
			\itemch  \textbf{Sai.} \\
			$\cos (a+b)=\cos a\cos  b-\sin a\sin b=\dfrac{15}{17}\cdot\dfrac{12}{13}-\dfrac{8}{17}\cdot\dfrac{5}{13}=\dfrac{140}{221}$ nên là mệnh đề sai.
			\itemch \textbf{Sai.} \\ 
			$\tan (a+b)=\dfrac{\tan  a+\tan b}{1-\tan a\tan b}=\dfrac{\dfrac{8}{15}+\dfrac{5}{12}}{1-\dfrac{8}{15}\cdot\dfrac{5}{12}}=\dfrac{171}{140}.$\\
			$\tan 2\left(a+b\right)=\dfrac{2\tan\left(a+b\right)}{1-\tan^2\left(a+b\right)}=\dfrac{2\cdot\dfrac{171}{140}}{1-\left(\dfrac{171}{140}\right)^2}\approx-5$  nên là mệnh đề sai.
		\end{itemchoice}
	 }
\end{ex}

% nguồn: Ok_dungsai_lop11_Lan2
\begin{ex}%[1D1H3-5]%[ma:f182de37]
	Cho $\tan a=\dfrac{1}{2},0 < a <\dfrac{\pi}{2};\tan b=\dfrac{1}{3},0 < b <\dfrac{\pi}{2}$. Xét tính đúng-sai của  các phát biểu sau
	 \choiceTF
		 {\True $\cot a=2$}
		{\True $\tan  (a+b)=1$}
		 {$\tan (a-b)=7$}
		{\True $\sin  (a-2b)=-\dfrac{2\sqrt 5}{25}$}
	\loigiai{ 
		\begin{itemchoice}
			\itemch  \textbf{Đúng.} \\
			Ta có $\tan a.\cot a=1\Rightarrow\cot a=\dfrac{1}{\tan a}=\dfrac{1}{\dfrac{1}{2}}=2$. 
			\itemch  \textbf{Đúng.} \\
			Ta có  $\tan (a+b)=\dfrac{\tan a+\tan b}{1-\tan a\tan b}=\dfrac{\dfrac{1}{2}+\dfrac{1}{3}}{1-\dfrac{1}{2}\cdot\dfrac{1}{3}}=1$.
			\itemch  \textbf{Sai.} \\
			Ta có  $\tan (a-b)=\dfrac{\tan a-\tan b}{1+\tan a\tan b}=\dfrac{\dfrac{1}{2}-\dfrac{1}{3}}{1+\dfrac{1}{2}\cdot\dfrac{1}{3}}=\dfrac{1}{7}$.
			\itemch  \textbf{Đúng.} \\
			Ta  có\\
			${{\cos  }^2}a=\dfrac{1}{1+{{\tan }^2}a}=\dfrac{1}{1+{{( \dfrac{1}{2})}^2}}=\dfrac{4}{5}\Rightarrow \hoac{& \cos a=\dfrac{2\sqrt {5}}{5}(\text{N}) \\ & \cos a=-\dfrac{2\sqrt {5}}{5}(\text{L}) }\Rightarrow \sin a=\tan a\cdot \cos a=\dfrac{\sqrt {5}}{5}$.\\
			${{\cos }^2}b=\dfrac{1}{1+{{\tan }^2}b}=\dfrac{1}{1+{{( \dfrac{1}{3})}^2}}=\dfrac{9}{10}\Rightarrow \hoac{& \cos b=\dfrac{3\sqrt {10}}{10}(\text{N}) \\ & \cos b=-\dfrac{3\sqrt {10}}{10}(\text{L})  }\Rightarrow \sin b=\tan b\cdot \cos b=\dfrac{\sqrt {10}}{10}$.\\
			 Suy ra
			$\sin ( a-2b)=\sin  a\cos 2b-\cos a\sin 2b=\sin a( 1-2{{\sin }^2}b)-\cos a\cdot 2\sin b\cos b=-\dfrac{2\sqrt {5}}{25}$.
		\end{itemchoice}
	 }
\end{ex}

% nguồn: Ok_dungsai_lop11_Lan2
\begin{ex}%[1D1H3-5]%[ma:ababe10f]
	Biết $\tan\alpha=2$. Khi  đó
	 \choiceTF
		{$\cot\alpha=-\dfrac{1}{2}$} 
		 {\True $\cos 2\alpha=-\dfrac{3}{5}$}
		{\True $\sin  2\alpha=\dfrac{4}{5}$}
		 {\True $\tan 2\alpha=-\dfrac{4}{3}$}
	 \loigiai{
		Ta  có
		\begin{itemize}
			\item $\cos 2\alpha=2\cos  ^2\alpha-1=\dfrac{2}{\tan^2\alpha+1}-1=\dfrac{2}{5}-1=-\dfrac{3}{5}$. \item $\sin 2\alpha=2\sin\alpha\cos\alpha=2\tan\alpha{\cos ^2}\alpha=\dfrac{2\tan\alpha}{\tan^2\alpha+1}=\dfrac{4}{5}$.
			\item $\tan 2\alpha=\dfrac{\sin 2\alpha}{\cos  2\alpha}=-\dfrac{4}{3}$.
		\end{itemize}
	 }
\end{ex}

% nguồn: Ok_dungsai_lop11_Lan2
\begin{ex}%[1D1H3-5]%[ma:a530994e]
	Cho  biết $\sin x=\dfrac{1}{\sqrt 3}$ và $ 0 < x <\dfrac{\pi}{2}$, khi đó
	 \choiceTF
		 {\True $\cos x > 0$}
		{\True $\cos x=\dfrac{\sqrt 6}{3}$} 
		{$\tan x=\dfrac{\sqrt  3}{3}$}
		 {$\cos\left(x+\dfrac{\pi}{3}\right)=\dfrac{\sqrt 6-3}{8}{\rm}$}
	\loigiai{ 
		Vì $ 0 < x <\dfrac{\pi}{2}$ nên $\cos  x > 0$. Ta có\\
		\begin{itemize}
			\item $\sin x=\dfrac{1}{\sqrt 3}$ $\Rightarrow $ $\cos x=\sqrt{1-\sin^2x}=\sqrt{1-\dfrac{1}{3}}=\dfrac{\sqrt 6}{3}$ $\Rightarrow $ $\tan x=\dfrac{\sin x}{\cos x}=\dfrac{\sqrt 2}{2}$. 
			\item $\cos\left(x+\dfrac{\pi}{3}\right)=\cos x\cos\dfrac{\pi}{3}-\sin x\sin\dfrac{\pi}{3}=\dfrac{\sqrt 6}{3}\cdot\dfrac{1}{2}-\dfrac{1}{\sqrt 3}\cdot\dfrac{\sqrt  3}{2}=\dfrac{\sqrt 6-3}{6}$.
		\end{itemize}
	} 
\end{ex}

% nguồn: Ok_dungsai_lop11_Lan2
\begin{ex}%[1D1V3-5]%[ma:f7636f07]
	 Xét tính đúng sai của các mệnh đề sau
	\choiceTF 
		 {$\sin ^2x=\dfrac{1+\sin 2x}{2}$}
		{\True Nếu  $\cos\alpha=\dfrac{1}{3}$ thì $\cos 2\alpha=-\dfrac{7}{9}$}
		{\True Nếu $\sin x=\dfrac{3}{4}$ với $ x\in\left(0;\dfrac{\pi}{2}\right)$ thì $\sin 2x=\dfrac{3\sqrt 7}{8}$} 
		{\True Cho $\cos\alpha=\dfrac{2}{3}$ với $\alpha\in\left(-\dfrac{\pi}{2};0\right)$ biết $\tan\left(\alpha+\dfrac{\pi}{4}\right)=a+b\sqrt c $. Khi  đó $ a+b+c=0$}
	 \loigiai{
	\begin{itemchoice}
		 \itemch Ta có $\sin ^2x=\dfrac{1-\cos 2x}{2}$.
		\itemch Ta có $\cos  2\alpha=2\cos ^2\alpha-1=2\left(\dfrac{1}{3}\right)^2-1=\dfrac{-7}{9}$.
		\itemch Ta có $\cos ^2x=1-\sin  ^2x=1-\left(\dfrac{3}{4}\right)^2=\dfrac{7}{16}$.\\
		Vì $ x\in\left(0;\dfrac{\pi}{2}\right)$  nên $\cos x > 0\Rightarrow\cos x=\dfrac{\sqrt 7}{4}$.\\
		Khi đó $\sin 2x=2\sin x\cdot\cos x=2\cdot\dfrac{\sqrt 7}{4}\cdot\dfrac{3}{4}=\dfrac{3\sqrt  7}{8}$.
		 \itemch Ta có $\tan ^2\alpha=\dfrac{1}{\cos^2\alpha}-1=\dfrac{1}{\left(\dfrac{2}{3}\right)^2}-1=\dfrac{5}{4}$\\
		Vì $\alpha\in\left(-\dfrac{\pi}{2};0\right)$ nên $\tan\alpha < 0\Rightarrow\tan\alpha=\dfrac{-\sqrt 5}{2}$.\\ 
		Khi đó $\tan\left(\alpha+\dfrac{\pi}{4}\right)=\dfrac{\tan\alpha+\tan\dfrac{\pi}{4}}{1-\tan\alpha\cdot\tan\dfrac{\pi}{4}}=\dfrac{\dfrac{-\sqrt 5}{2}+1}{1-\left(\dfrac{-\sqrt 5}{2}\right)\cdot1}=-9+4\sqrt 5$.\\ 
		Vậy $  a=-9, b=4, c=5$ nên suy ra $a+b+c=-9+4+5=0$.
	\end{itemchoice}
	} 
\end{ex}

% nguồn: Ok_dungsai_lop11_Lan2
\begin{ex}%[1D1V3-5]%[ma:91a1a8f8]
	Xét tính  đúng sai của các phát biểu sau
	 \choiceTF
		{\True $\sin\left(a+b\right)=\sin a\cos b+\cos  a\sin b$}
		{$\sin{31^\circ}\cdot\cos{12^\circ}+\cos{12^\circ}\cdot\sin{31^\circ}=\sin{19^\circ}$} 
		{\True  Cho $\cos x=\dfrac{4}{5},x\in\left(-\dfrac{\pi}{2};0\right)$. Giá trị của $\sin 2x$ là $-\dfrac{24}{25}$}
		{Cho $\sin\alpha=\dfrac{1}{3}$ và $\dfrac{\pi}{2}<\alpha  <\pi$. Biết giá trị của $\cos\left(\alpha-\dfrac{\pi}{6}\right)=\dfrac{1-a\sqrt 6}{b}$ với $a, b\in\mathbb{N}$ thì $a+b=4$}
	\loigiai{ 
	\begin{itemchoice}
		 \itemch Theo công thức cộng, ta có $\sin\left(a+b\right)=\sin a\cos b+\cos a\sin b$.
		\itemch Áp dụng công thức  cộng, ta có $\sin{31^{\circ}}\cdot\cos{12^\circ}+\cos{12^\circ}\cdot\sin{31^\circ}=\sin\left(31^\circ+12^\circ\right)=\sin{43^\circ}$.
		\itemch Ta có $\sin ^2x=1-\cos ^2x=1-\dfrac{16}{25}=\dfrac{9}{25}$. Mà $x\in\left(-\dfrac{\pi}{2};0\right)\Rightarrow\sin x=-\dfrac{3}{5}$.\\ 
		Vậy $\sin 2x=2\sin x\cdot\cos x=2\cdot\dfrac{4}{5}\cdot\left(-\dfrac{3}{5}\right)=-\dfrac{24}{25}$. 
		\itemch Vì $\sin\alpha=\dfrac{1}{3},\dfrac{\pi}{2}<\alpha <\pi$ nên $\cos\alpha=-\dfrac{2\sqrt 2}{3}$.\\ 
		Do đó $\cos\left(\alpha-\dfrac{\pi}{6}\right)=\cos\alpha  \cdot\cos\dfrac{\pi}{6}+\sin\alpha \cdot\sin\dfrac{\pi}{6}=-\dfrac{2\sqrt 2}{3}\cdot\dfrac{\sqrt 3}{2}+\dfrac{1}{3}\cdot\dfrac{1}{2}=\dfrac{1-2\sqrt 6}{6}$.\\
		Khi đó $ a=2,\,\,b=6\Rightarrow  a+b=8$.
	\end{itemchoice}
	 }
\end{ex}

% nguồn: Ok_dungsai_lop11_Lan2
\begin{ex}%[1D1V3-5]%[ma:b1836fe6]
	 Cho $\cos a=\dfrac{3}{4}$, $\sin a > 0$, $\sin b=\dfrac{3}{5}$, $\cos b < 0$ . Khi đó
	 \choiceTF
		{\True Giá trị của $\tan  a=\dfrac{\sqrt 7}{3}$}
		{Giá trị của $\cot  b=-\dfrac{2}{3}$}
		{Giá trị của $\cos 2a+\cos 2b$ thuộc  khoảng $\left(\dfrac{1}{2};1\right)$}
		{Giá trị của $\cos\left(a+b\right)$ thuộc khoảng $\left(-\dfrac{1}{2};-\dfrac{1}{3}\right)$} 
	\loigiai{ 
	\begin{itemchoice}
		\itemch Ta có $\tan  ^2a+1=\dfrac{1}{\cos^2a}\Rightarrow\tan a=\sqrt{\dfrac{1}{\cos^2a}-1}=\sqrt{\dfrac{16}{9}-1}=\dfrac{\sqrt 7}{3}$.
		\itemch Ta có $\cot ^2b+1=\dfrac{1}{\sin^2b}\Rightarrow\cot  b=-\sqrt{\dfrac{1}{\sin^2a}-1}=-\sqrt{\dfrac{25}{9}-1}=-\dfrac{4}{3}$.
		\itemch Ta có $\cos 2a+\cos  2b=2\cos ^2a-1+1-2\sin ^2b=2\cdot\dfrac{9}{16}-2\cdot\dfrac{9}{25}=\dfrac{81}{200}=0{,}405\notin\left(\dfrac{1}{2};1\right)$.
		 \itemch Ta có
		\begin{itemize}
			\item $\left\{\begin{array}{l}
			\cos  a=\dfrac{3}{4}\\
			\sin a > 0 
		\end{array}\right.\Rightarrow\sin a=\sqrt{1-\cos^2a}=\dfrac{\sqrt 7}{4}$.
		\item $\left\{\begin{array}{l}
		 \sin b=\dfrac{3}{5}\\
		\cos b < 0 
	\end{array}\right.\Rightarrow\cos b=-\sqrt{1-\sin^2b}=-\dfrac{4}{5}$.
\end{itemize}
Suy  ra
\begin{align*}
	 \cos\left(a+b\right)&\,=\cos a\cos b-\sin a\sin b\\
	&\,=\dfrac{3}{4}\cdot\left(-\dfrac{4}{5}\right)-\dfrac{\sqrt  7}{4}\cdot\dfrac{3}{5}\\
	&\,=-\dfrac{3}{5}\left(1+\dfrac{\sqrt 7}{4}\right)\notin\left(-\dfrac{1}{2};-\dfrac{1}{3}\right). 
\end{align*}
\end{itemchoice}
} 
\end{ex}

% nguồn: Ok_dungsai_lop11_Lan2
\begin{ex}%[1D1V3-5]%[ma:caf8b8b4]
	Xét tính đúng sai của các  mệnh đề sau
	 \choiceTF
		{$\sin\left(x-\dfrac{\pi}{6}\right)=\dfrac{\sqrt 3}{2}\sin x+\dfrac{1}{2}\cos  x$}
		{\True $4\sin x\cdot\sin 2x\cdot\sin 3x=\sin  4x+\sin 2x-\sin 6x$}
		{$ 1+\sin 2x+\cos 2x=2\sqrt 2\sin x\cdot  \cos\left(x-\dfrac{\pi}{4}\right)$}
		{$\sin ^2x+2\sin\left(a-x\right)\cdot\sin  x\cdot\cos a+\sin ^2\left(a-x\right)=\cos ^2a$}
	 \loigiai{
	\begin{itemchoice}
		\itemch Theo công thức cộng ta có $$\sin\left(x-\dfrac{\pi}{6}\right)=\sin x\cdot\cos\dfrac{\pi}{6}-\cos x\cdot\sin\dfrac{\pi}{6}=\dfrac{\sqrt 3}{2}\sin x-\dfrac{1}{2}\cos  x.$$
		\itemch Áp dụng công  thức biến đổi từ tích thành tổng, ta được
		$$4\sin x\cdot\sin 2x\cdot\sin 3x=2\left(\cos x-\cos 3x\right)\sin 3x=\sin  4x+\sin 2x-\sin 6x.$$
		 Do đó $4\sin x\cdot\sin 2x\cdot\sin 3x=\sin 4x+\sin 2x-\sin 6x$.
		\itemch Ta có $ 1+\sin 2x+\cos 2x=2\sin x\cos x+2\cos ^2x=2\cos  x(\sin x+\cos x)$.\\
		Mà $\sin x+\cos x=\sqrt 2\left(\cos x\dfrac{1}{\sqrt 2}+\sin x\dfrac{1}{\sqrt 2}\right)$ nên suy ra $2\cos x(\sin x+\cos x)=\sqrt 2\cos\left(x-\dfrac{\pi}{4}\right)$. 
		\itemch  Ta có $\sin ^2x+2\sin\left(a-x\right)\cdot\sin x\cdot\cos a+\sin ^2\left(a-x\right)$\\
		$\begin{array}{l}
		=\sin  ^2x+\sin\left(a-x\right)\cdot\left[\sin\left(x+a\right)+\sin\left(x-a\right)\right]+\sin ^2\left(a-x\right)\\
		=\sin ^2x+\sin\left(a-x\right)\cdot\left[\sin\left(x+a\right)+\sin\left(x-a\right)+\sin\left(a-x\right)\right]\\ 
		=\sin ^2x+\sin\left(a-x\right)\cdot\left[\sin\left(x+a\right)+\sin\left(x-a\right)-\sin\left(x-a\right)\right]\\ 
		 =\sin ^2x+\sin\left(a-x\right)\cdot\sin\left(x+a\right)\\
		 =\dfrac{1-\cos 2x}{2}-\dfrac{\cos 2a-\cos 2x}{2}\\
		=\dfrac{1-\cos  2a}{2}=\sin ^2a.
	\end{array}$\\
\end{itemchoice}
} 
\end{ex}

% nguồn: Ok_dungsai_lop11_Lan2
\begin{ex}%[1D1V3-5]%[ma:0b91b7a1]
	Biết $0<a, b<\dfrac{\pi}{2}$,  $a+b=\dfrac{\pi}{4}$ và $\tan{a}\tan{b}=3-2\sqrt{2}$.
	 \choiceTF
		{\True $\tan{a}+\tan{b}=-2+2\sqrt{2}$} 
		{\True $\tan{a}=-1+\sqrt{2}$} 
		 {$\tan{b}=-1-\sqrt{2}$}
		 {$\tan{a}-\tan{b}=-2-2\sqrt{2}$}
	\loigiai{ 
	Ta có $\tan{(a+b)}=\dfrac{\tan{a}+\tan{b}}{1-\tan{a}\tan{b}}\Rightarrow  \tan{a}+\tan{b}=\tan{(a+b)}(1-\tan{a}\tan{b})$.\\
	Mà $a+b=\dfrac{\pi}{4}$  và $\tan{a}\tan{b}=3-2\sqrt{2}$ nên suy ra \[\tan{a}+\tan{b}=\tan{\dfrac{\pi}{4}}\left[1-\left(3-2\sqrt{2}\right)\right]=-2+2\sqrt{2}.\]
	Đặt $S=\tan{a}+\tan{b}; P=\tan{a}\tan{b}$. Khi đó  $\tan{a}$ và $\tan{b}$ là nghiệm của phương trình
	 \[X^2-SX+P=0\Leftrightarrow X^2-\left(2\sqrt{2}-2\right)X+3-2\sqrt{2}=0\Leftrightarrow X=-1+\sqrt{2}.\]
	Suy ra $\tan{a}=\tan{b}=-1+\sqrt{2}$.\\ 
	Khi đó  $\tan{a}-\tan{b}=0$.
	} 
\end{ex}

% nguồn: Ok_dungsai_lop11_Lan2
\begin{ex}%[1D1V3-5]%[ma:cf54db80]
	Cho góc $\alpha$ thỏa  mãn $-\dfrac{\pi}{2}<\alpha<0$ và $\cot{\alpha}=-3$.
	\choiceTF 
		{$\sin{\alpha}>0$} 
		 {\True $\sin{\alpha}=-\dfrac{\sqrt{10}}{10}$}
		{\True  $\dfrac{\cos{\alpha}-\sin{\alpha}}{\cos^3{\alpha}+3\sin^3{\alpha}+2\cos{\alpha}}=\dfrac{10}{21}$}
		{\True $\left[\tan{\dfrac{17\pi}{4}}+\tan{\left(\dfrac{7\pi}{2}-\alpha\right)}\right]^2+\left[\cot{\dfrac{13\pi}{4}}+\cot{(7\pi-\alpha)}\right]^2=20$} 
	\loigiai{ 
	Ta có  $1+\cot^2{\alpha}=\dfrac{1}{\sin^2{\alpha}}\Rightarrow\sin^2{\alpha}=\dfrac{1}{1+\cot^2{\alpha}}=\dfrac{1}{1+(-3)^2}=\dfrac{1}{10}$.\\
	Mà $-\dfrac{\pi}{2}<\alpha<0$  nên suy ra $\sin{\alpha}=-\dfrac{\sqrt{10}}{10}$.\\
	 Đặt $P=\dfrac{\cos{\alpha}-\sin{\alpha}}{\cos^3{\alpha}+3\sin^3{\alpha}+2\cos{\alpha}}$.\\
	Chia cả tử và mẫu của $P$ cho  $\sin^3{\alpha}$, ta được
	\[P=\dfrac{\dfrac{\cos{\alpha}}{\sin^3{\alpha}}-\dfrac{1}{\sin^2{\alpha}}}{\dfrac{\cos^3{\alpha}}{\sin^3{\alpha}}+3+2\dfrac{\cos{\alpha}}{\sin^3{\alpha}}}=\dfrac{\cot{\alpha}\left(1+\cot^2{\alpha}\right)-\left(1+\cot^2{\alpha}\right)}{\cot^3{\alpha}+3+2\cot{\alpha}\left(1+\cot^2{\alpha}\right)}=\dfrac{10}{21}.\] 
	Ta có 
	\begin{itemize}
		\item $\tan{\dfrac{17\pi}{4}}=\tan{\left(\dfrac{\pi}{4}+4\pi\right)}=\tan{\dfrac{\pi}{4}}=1;  \tan{\left(\dfrac{7\pi}{2}-\alpha\right)}=\cot{\alpha}$.
		 \item $\cot{\left(\dfrac{13\pi}{4}\right)}=\cot{\left(\dfrac{\pi}{4}+3\pi\right)}=\cot{\dfrac{\pi}{4}}=1; \cot{(7\pi-\alpha)}=-\cot{\alpha}$.
	\end{itemize}
	Khi đó 
	\begin{align*}
		&\,\left[\tan{\dfrac{17\pi}{4}}+\tan{\left(\dfrac{7\pi}{2}-\alpha\right)}\right]^2+\left[\cot{\dfrac{13\pi}{4}}+\cot{(7\pi-\alpha)}\right]^2\\ 
		 &\,=\left(1+\cot{\alpha}\right)^2+\left(1-\cot{\alpha}\right)^2\\
		 &\,=2+2\cot^2{\alpha}=2+2\cdot(-3)^2=20.
	\end{align*}
	} 
\end{ex}


%%%%%%%%%%%%%%%% TRẢ LỜI NGẮN %%%%%%%%%%%%%%%%

% nguồn: Ok_TraLoiNgan_Lop11_Lan2
\begin{ex}%[1D1H3-5]%[ma:8ed4442f]
	Rút gọn $C=\tan 3x-\tan x-\dfrac{2\sin x}{\cos 3x}$. 
	 \shortans{0}
	\loigiai{ 
		 Ta có
		\allowdisplaybreaks 
		\begin{eqnarray*}
			C&=&\tan 3x-\tan  x-\dfrac{2\sin x}{\cos 3x}=\dfrac{\sin 3x}{\cos 3x}-\dfrac{\sin x}{\cos x}-\dfrac{2\sin x}{\cos 3x}\\
			&=&\dfrac{\sin 3x\cos x-\sin  x\cos 3x}{\cos 3x\cos x}-\dfrac{2\sin x}{\cos 3x}=\dfrac{\sin 2x}{\cos 3x\cos x}-\dfrac{2\sin x}{\cos 3x}\\
			&=&\dfrac{2\sin  x\cos x}{\cos 3x\cos x}-\dfrac{2\sin x}{\cos 3x}=\dfrac{2\sin x}{\cos 3x}-\dfrac{2\sin x}{\cos 3x}=0.
		\end{eqnarray*}
		 Do đó $C=0$.
	} 
\end{ex}

% nguồn: Ok_TraLoiNgan_Lop11_Lan2
\begin{ex}%[1D1H3-5]%[ma:626f74b6]
	Cho biểu thức $g(x)=\dfrac{-2\cot ^2x-\cos ^2x}{\cot ^2x}+\dfrac{\sin  x \cos x}{\cot x}$ với $x \neq 0$, $x \neq \dfrac{\pi}{2}$, $x \neq \pi$.
	Tính $g\left(\dfrac{2024\pi}{2023}\right)$. 
	 \shortans{-2}
	 \loigiai{
		Ta có 
		\begin{eqnarray*}
			g(x) &=&-2-\dfrac{\cos^2 x}{\cot^2 x}+\dfrac{\sin x  \cos x}{\dfrac{\cos x}{\sin x}} \\
			 &=&-2-\dfrac{\cos^2 x}{\dfrac{\cos^2 x}{\sin^2 x}}+\sin^2 x\\
			&=&-2-\sin^2 x+\sin^2  x\\
			&=&-2. 
		\end{eqnarray*}
	 }
\end{ex}

% nguồn: Ok_TraLoiNgan_Lop11_Lan2
\begin{ex}%[1D1H3-5]%[ma:e767044d]
	Cho $\alpha+\beta=\dfrac{\pi}{3}$ và $\cos\alpha\ne\cos\beta$. Tính $\dfrac{\sin\alpha-\sin\beta}{\cos\alpha-\cos\beta}=a$ (làm tròn kết quả  đến hàng phần mười).\\
	 \shortans{-1,7}
	 \loigiai{
	$\dfrac{\sin\alpha-\sin\beta}{\cos\alpha-\cos\beta}=\dfrac{2\cos\dfrac{\alpha+\beta}{2}\sin\dfrac{\alpha-\beta}{2}}{-2\sin\dfrac{\alpha+\beta}{2}\sin\dfrac{\alpha-\beta}{2}}=-\cot\dfrac{\alpha+\beta}{2}=-\cot\dfrac{\pi}{6}=-\sqrt 3 \approx-1{,7}$.} 
\end{ex}

% nguồn: Ok_TraLoiNgan_Lop11_Lan2
\begin{ex}%[1D1H3-5]%[ma:bb4d5a54]
	 Cho $\cot\alpha=3$. Tính $A=\dfrac{3\sin\alpha-2\cos\alpha}{12\sin^3\alpha+4\cos^3\alpha}$.
	 \shortans{0,25}
	 \loigiai{
		 Ta có
		 $A=\dfrac{3\sin\alpha-2\cos\alpha}{12\sin^3\alpha+4\cos^3\alpha}=\dfrac{\dfrac{1}{\sin^2\alpha}(3-2\cot\alpha)}{12+4\cot^3\alpha}=\left(1+\cot^2\alpha\right)\dfrac{3-2\cot\alpha}{12+4\cot^3\alpha}=-\dfrac{1}{4}$.
	} 
\end{ex}

% nguồn: Ok_TraLoiNgan_Lop11_Lan2
\begin{ex}%[1D1H3-5]%[ma:35b97801]
	Cho hai góc nhọn $a$ và $b$. Biết $\cos a=\dfrac{1}{3}$  và $\cos b=\dfrac{1}{4}$. Tính giá trị của \[P=(\cos a \cos b)^2-(\sin a \sin b)^2.\]
	 \shortans{-0,8}
	\loigiai{ 
		 Ta có
		\[ 
		\begin{aligned}
			& P=(\cos a \cos b)^2-(\sin a \sin  b)^2=(\cos a \cos b)^2-\left(1-\cos^2 a\right)\left(1-\cos^2 b\right) \\
			& \quad=\left(\dfrac{1}{12}\right)^2-\dfrac{8}{9} \cdot \dfrac{15}{16}=-\dfrac{119}{144}\approx  -0{,}8. \\
		\end{aligned}
		 \]
	} 
\end{ex}

% nguồn: Ok_TraLoiNgan_Lop11_Lan2
\begin{ex}%[1D1H3-5]%[ma:d915def2]
	Cho $\tan\alpha=2$ .  Tính $C=\dfrac{\sin\alpha}{\sin^3\alpha+2\cos^3\alpha}$.
	\shortans{1} 
	\loigiai{ 
		Tacó  $C=\dfrac{\sin\alpha}{\sin^3\alpha+2\cos^3\alpha}=\dfrac{\tan\alpha\cdot\dfrac{1}{\cos^2\alpha}}{\tan^3\alpha+2}=\dfrac{\tan\alpha\left(1+\tan^2\alpha\right)}{\tan^3\alpha+2}=\dfrac{2\left(1+2^2\right)}{2^3+2}=1$.
	 }
\end{ex}

% nguồn: Ok_TraLoiNgan_Lop11_Lan2
\begin{ex}%[1D1H3-5]%[ma:9c8e51a1]
	 Biết $\dfrac{3-4\cos 2a+\cos 4a}{3+4\cos 2a+\cos 4a}=k{\tan ^4}a$. Tìm $k$?
	 \shortans{1}
	\loigiai{ 
		Ta  có
		\begin{eqnarray*}
			\dfrac{3-4\cos 2a+\cos 4a}{3+4\cos 2a+\cos 4a}&=&\dfrac{2\cos^22a-4\cos  2a+2}{2\cos^22a+4\cos 2a+2}\\
			 &=&\dfrac{2\left(\cos^22a-2\cos 2a+1\right)}{2\left(\cos^22a+2\cos 2a+1\right)}\\
			 &=&\dfrac{(1-\cos 2a)^2}{(1+\cos 2a)^2}\\
			&=&\dfrac{\left[1-\left(1-2\sin^2a\right)\right]^2}{\left[1+\left(2\cos^2a-1\right)\right]^2}\\ 
			 &=&\dfrac{4\sin^4a}{4\cos^4a}=\tan ^4a.
		\end{eqnarray*}
	} 
\end{ex}

% nguồn: Ok_TraLoiNgan_Lop11_Lan2
\begin{ex}%[1D1H3-5]%[ma:36ae66ad]
	 Rút gọn biểu thức $\sin 6a\cot 3a-\cos 6a$.
	 \shortans{1}
	\loigiai{ 
		Ta  có
		\begin{eqnarray*}
			\sin  6a\cot 3a-\cos 6a&=&2\sin 3a\cos 3a\dfrac{\cos 3a}{\sin 3a}-\left(2\cos^23a-1\right)\\
			&=&2\cos  ^23a-\left(2\cos^23a-1\right)\\
			&=&1\, (\text{không phụ thuộc} \,a). 
	\end{eqnarray*}}
\end{ex}

% nguồn: Ok_TraLoiNgan_Lop11_Lan2
\begin{ex}%[1D1H3-5]%[ma:77ded8e7]
	Cho tam giác $ABC$ có $AB=c,AC=b,AB=c$, thỏa  $\dfrac{-1+\cos (\pi+B)}{\sin (\pi-A-C)}=-\dfrac{2a+c}{\sqrt{4a^2-c^2}}$. Khi đó tam giác $ABC$ có $ma=nb$. Tính tổng $m+n$.
	\shortans{2} 
	 \loigiai{Ta có
		\begin{eqnarray*}
			&&\dfrac{-1+\cos  (\pi+B)}{\sin (\pi-A-C)}=-\dfrac{2a+c}{\sqrt{4a^2-c^2}}\\
			&\Leftrightarrow&\dfrac{-1-\cos B}{\sin  (\pi-(A+C))}=-\dfrac{2a+c}{\sqrt{4a^2-c^2}}\\
			&\Leftrightarrow&\dfrac{1+\cos B}{\sin B}=\dfrac{2a+c}{\sqrt{4a^2-c^2}}\\ 
			&\Leftrightarrow&\dfrac{1+\cos B}{\sin  B}=\dfrac{2a+c}{\sqrt{4a^2-c^2}}\\
			&\Leftrightarrow&{\left(\dfrac{1+\cos B}{\sin B}\right)^2}=\left(\dfrac{2a+c}{\sqrt{4a^2-c^2}}\right)^2\\ 
			&\Leftrightarrow&\dfrac{(1+\cos  B)^2}{\sin^2B}=\dfrac{(2a+c)^2}{4a^2-c^2}\\
			 &\Leftrightarrow&\dfrac{(1+\cos B)^2}{1-\cos^2B}=\dfrac{(2a+c)^2}{4a^2-c^2}\\
			&\Leftrightarrow&\dfrac{1+\cos  B}{1-\cos B}=\dfrac{2a+c}{2a-c}\\
			 &\Leftrightarrow&\dfrac{1+\cos B}{1-\cos B}-1=\dfrac{2a+c}{2a-c}-1\\
			&\Leftrightarrow&2ac\cdot\cos B=c^2\\ 
			 &\Leftrightarrow&{a^2}+c^2-b^2=c^2\\
			&\Leftrightarrow&{a^2}=b^2\Leftrightarrow  a=b.
		\end{eqnarray*}
		 Khi đó $m+n=2$.
	} 
\end{ex}

% nguồn: Ok_TraLoiNgan_Lop11_Lan2
\begin{ex}%[1D1H3-5]%[ma:007e56e5]
	 Rút gọn biểu thức $(\tan a-\tan b)\cot (a-b)-\tan a\tan b$.
	 \shortans{1}
	 \loigiai{
		 Ta có
		\begin{eqnarray*}
			(\tan a-\tan b)\cot (a-b)-\tan a\tan b&=&\dfrac{\tan a-\tan b}{\tan  (a-b)}-\tan a\tan b\\
			&=&\dfrac{\tan a-\tan  b}{\dfrac{\tan a-\tan b}{1+\tan a\tan b}}-\tan a\tan b\\
			&=&1+\tan a\tan b-\tan  a\tan b\\
			&=&1\, (\text{không phụ thuộc}\,  a).
		\end{eqnarray*}
	} 
\end{ex}

% nguồn: Ok_TraLoiNgan_Lop11_Lan2
\begin{ex}%[1D1V3-5]%[ma:51178a72]
	Cho $\tan x=-\dfrac{4}{3}$ và $\dfrac{\pi}{2} < x < \pi$. Tính giá trị của biểu  thức $M=\dfrac{\sin^2 x-\cos x}{\sin x-\cos^2 x}$. (Kết quả làm tròn đến hàng phần mười).
	\shortans{2,8} 
	 \loigiai{
		Ta có $\tan  x=-\dfrac{4}{3} \Rightarrow \cos ^2x=\dfrac{1}{1+\tan ^2x}=\dfrac{9}{25} \Rightarrow \cos x=\pm \dfrac{3}{5}$.\\
		Vì $\dfrac{\pi}{2} < x < \pi \Rightarrow \cos  x=-\dfrac{3}{5} \Rightarrow \sin x=\tan x \cdot \cos x=\dfrac{4}{5}$. \\
		Khi đó $M=\dfrac{\sin ^2x-\cos x}{\sin x-\cos ^2x}=\dfrac{31}{11}\approx 2{,}8$. 
	} 
\end{ex}

% nguồn: Ok_TraLoiNgan_Lop11_Lan2
\begin{ex}%[1D1V3-5]%[ma:9d79ad40]
	Tính $N=\sin ^2\alpha \cdot \tan  ^2\alpha+4\sin ^2\alpha-\tan ^2\alpha+3\cos ^2\alpha$.
	 \shortans{3}
	 \loigiai{
		Ta  có
		\[ 
		\begin{aligned}
			N  &=\sin^2 \alpha \cdot \tan^2 \alpha+4 \sin^2 \alpha-\tan^2 \alpha+3 \cos^2 \alpha=\tan^2 \alpha\left(\sin^2 \alpha-1\right)+4 \sin^2 \alpha+3 \cos^2 \alpha \\
			&=-\tan^2 \alpha \cos^2 \alpha+4 \sin^2 \alpha+3 \cos^2 \alpha=-\sin^2 \alpha+4 \sin^2 \alpha+3 \cos^2  \alpha=3.
		\end{aligned}
		\] 
	} 
\end{ex}

% nguồn: Ok_TraLoiNgan_Lop11_Lan2
\begin{ex}%[1D1V3-5]%[ma:11fed84f]
	Tính $M=\sqrt{\sin ^4\alpha+4\cos ^2\alpha}+\sqrt{\cos ^4\alpha+4\sin ^2\alpha}$ khi $\alpha=\dfrac{8}{7} \pi$. 
	 \shortans{3}
	 \loigiai{
		Ta có 
		 \[
		\begin{aligned}
			M &=\sqrt{\sin^4 \alpha+4 \cos^2  \alpha}+\sqrt{\cos^4 \alpha+4 \sin^2 \alpha} \\
			&=\sqrt{\left(2-\sin^2  \alpha\right)^2}+\sqrt{\left(2-\cos^2 \alpha\right)^2}=\left|2-\sin^2 \alpha\right|+\left|2-\cos^2 \alpha\right|\\
			 &=2-\sin^2 \alpha+2-\cos^2 \alpha=4-1=3.
		\end{aligned}
		\] 
	 }
\end{ex}

% nguồn: Ok_TraLoiNgan_Lop11_Lan2
\begin{ex}%[1D1V3-5]%[ma:ab63a905]
	Cho  biểu thức $f(x)=3\left(\sin^4x+\cos ^4x\right)-2\left(\sin ^6x+\cos ^6x\right)$ tính $f(1)$.
	\shortans{1} 
	 \loigiai{
		Ta có $\sin ^4x+\cos ^4x=1-2\sin ^2x \cos ^2x,  \sin ^6x+\cos ^6x=1-3\sin ^2x \cos ^2x$.\\
		Suy ra $f(x)=3\left(1-2\sin ^2x \cos  ^2x\right)-2\left(1-3\sin ^2x \cos ^2x\right)=1$.
	} 
\end{ex}

% nguồn: Ok_TraLoiNgan_Lop11_Lan2
\begin{ex}%[1D1V3-5]%[ma:d13bac88]
	Tính giá trị của biểu thức  $M=\sin\dfrac{\pi}{24}\sin\dfrac{5\pi}{24}\sin\dfrac{7\pi}{24}\sin\dfrac{11\pi}{24}$ (làm tròn đến hàng phần trăm).
	 \shortans{0,06}
	\loigiai{ 
	Ta có $\sin\dfrac{7\pi}{24}=\cos\dfrac{5\pi}{24}$ và $\sin\dfrac{11\pi}{24}=\cos\dfrac{\pi}{24}$.\\ 
	Do đó 
	\begin{align*}
		 M=&\,\sin\dfrac{\pi}{24}\sin\dfrac{5\pi}{24}\cos\dfrac{5\pi}{24}\cos\dfrac{\pi}{24}\\
		 =&\,\dfrac{1}{4}\cdot\left(2\cdot\sin\dfrac{\pi}{24}\cdot\cos\dfrac{\pi}{24}\right)\cdot\left(2\cdot\sin\dfrac{5\pi}{24}\cdot\cos\dfrac{5\pi}{24}\right)\\
		=&\,\dfrac{1}{4}\cdot\sin\dfrac{\pi}{12}\cdot\sin\dfrac{5\pi}{12}=\dfrac{1}{4}\cdot\dfrac{1}{2}\left(\cos\dfrac{6\pi}{12}+\cos\dfrac{\pi}{3}\right)\\ 
		 =&\,\dfrac{1}{8}\cdot\left(0+\dfrac{1}{2}\right)=\dfrac{1}{16}\approx 0{,}06.
	\end{align*}
	} 
\end{ex}

% nguồn: Ok_TraLoiNgan_Lop11_Lan2
\begin{ex}%[1D1V3-5]%[ma:bca6e47d]
	Cho góc  $\alpha $ thỏa mãn $\dfrac{\pi}{2}<\alpha <\pi$ và $\sin\alpha=\dfrac{4}{5}$. Tính $P=\sin 2\left(\alpha+\pi\right)$ (làm tròn đến hàng đơn vị).
	 \shortans{-1}
	 \loigiai{
	Ta có  $P=\sin 2\left(\alpha+\pi\right)=\sin\left(2\alpha+2\pi\right)=\sin 2\alpha=2\sin\alpha\cos\alpha $ .\\
	Từ hệ thức $\sin ^2\alpha+\cos ^2\alpha=1$ ,  suy ra $\cos\alpha=\pm\sqrt{1-\sin^2\alpha}=\pm\dfrac{3}{5}$ .\\
	Do $\dfrac{\pi}{2}<\alpha  <\pi $ nên ta chọn $\cos\alpha=-\dfrac{3}{5}$ .\\
	Thay $\sin\alpha=\dfrac{4}{5}$ và $\cos\alpha=-\dfrac{3}{5}$ vào $ P$,  ta được $ P=2\cdot\dfrac{4}{5}\cdot\left(-\dfrac{3}{5}\right)=-\dfrac{24}{25}\approx-0{,}96$.
	} 
\end{ex}

% nguồn: Ok_TraLoiNgan_Lop11_Lan2
\begin{ex}%[1D1V3-5]%[ma:569c670c]
	Tính giá trị của biểu thức $\cos\dfrac{\pi}{30}\cos\dfrac{\pi}{5}+\sin\dfrac{\pi}{30}\sin\dfrac{\pi}{5}$ (làm tròn đến hàng  phần trăm).
	 \shortans{0,87}
	 \loigiai{
	Ta có $$\cos\dfrac{\pi}{30}\cos\dfrac{\pi}{5}+\sin\dfrac{\pi}{30}\sin\dfrac{\pi}{5}=\cos\left(\dfrac{\pi}{30}-\dfrac{\pi}{5}\right)=\cos\left(-\dfrac{\pi}{6}\right)=\dfrac{\sqrt 3}{2}\approx 0{,}87.$$ 
	 }
\end{ex}

% nguồn: Ok_TraLoiNgan_Lop11_Lan2
\begin{ex}%[1D1V3-5]%[ma:aee99121]
	Tính giá trị của biểu  thức $ M=\cos\dfrac{2\pi}{7}+\cos\dfrac{4\pi}{7}+\cos\dfrac{6\pi}{7}$.
	\shortans{-0,5} 
	\loigiai{ 
	Áp dụng công thức $\sin a-\sin  b=2\cdot\cos\dfrac{a+b}{2}\cdot\sin\dfrac{a-b}{2}.$\\
	 Ta có
	\begin{align*}
		2\sin\dfrac{\pi}{7}\cdot M=&\,2\cdot\cos\dfrac{2\pi}{7}\cdot\sin\dfrac{\pi}{7}+2\cdot\cos\dfrac{4\pi}{7}\cdot\sin\dfrac{\pi}{7}+2\cdot\cos\dfrac{6\pi}{7}\cdot\sin\dfrac{\pi}{7}\\ 
		=&\,\sin\dfrac{3\pi}{7}-\sin\dfrac{\pi}{7}+\sin\dfrac{5\pi}{7}-\sin\dfrac{3\pi}{7}+\sin\dfrac{7\pi}{7}-\sin\dfrac{5\pi}{7}\\ =&\,-\sin\dfrac{\pi}{7}+\sin\pi=-\sin\dfrac{\pi}{7}. 
	\end{align*}
	 Khi đó $M=-\dfrac{1}{2}$.
	 }
\end{ex}

% nguồn: Ok_TraLoiNgan_Lop11_Lan2
\begin{ex}%[1D1V3-5]%[ma:bbe247af]
	Biết $\sin\left(\pi-\alpha\right)=-\dfrac{3}{5}$ và $\pi <\alpha <\dfrac{3\pi}{2}$. Tính $ P=\sin\left(\alpha+\dfrac{\pi}{6}\right)$ (làm tròn đến hàng  phần chục).
	 \shortans{-0,9}
	 \loigiai{
	Ta có $-\dfrac{3}{5}=\sin\left(\pi-\alpha\right)=\sin\alpha $.\\ 
	Từ hệ thức $\sin ^2\alpha+\cos ^2\alpha=1$,  suy ra $\cos\alpha=\pm\sqrt{1-\sin^2\alpha}=\pm\dfrac{4}{5}$.\\
	 Do $\pi <\alpha <\dfrac{3\pi}{2}$ nên ta chọn $\cos\alpha=-\dfrac{4}{5}$.\\
	Suy ra $ P=\sin\left(\alpha+\dfrac{\pi}{6}\right)=\dfrac{\sqrt 3}{2}\sin\alpha+\dfrac{1}{2}\cos\alpha=\dfrac{\sqrt  3}{2}\left(-\dfrac{3}{5}\right)+\dfrac{1}{2}\left(-\dfrac{4}{5}\right)=\dfrac{-4-3\sqrt 3}{10}\approx-0,9$.
	 }
\end{ex}

% nguồn: Ok_TraLoiNgan_Lop11_Lan2
\begin{ex}%[1D1V3-5]%[ma:af76ea8f]
	Cho $\sin \alpha-\cos \alpha=\dfrac{1}{\sqrt{5}}$. Tính giá trị của $P=\sqrt{\sin  ^4\alpha+\cos ^4\alpha}$. (Kết quả làm tròn đến hàng phần trăm).
	\shortans{0,82} 
	\loigiai{ 
		 Ta có
		 \allowdisplaybreaks
		\begin{eqnarray*}
			&&\sin \alpha-\cos \alpha=\dfrac{1}{\sqrt{5}} \\ 
			&\Rightarrow&  (\sin \alpha-\cos \alpha)^2=\dfrac{1}{5}\\
			&\Leftrightarrow& 1-2\sin \alpha \cos \alpha=\dfrac{1}{5}\\  &\Leftrightarrow& \sin \alpha \cos \alpha=\dfrac{2}{5}.
		\end{eqnarray*}
		 Ta có
		\begin{eqnarray*}
			P&=&\sqrt{\sin ^4\alpha+\cos ^4\alpha}=\sqrt{\left(\sin ^2\alpha+\cos  ^2\alpha\right)^2-2\sin ^2\alpha \cos ^2\alpha}\\
			&=&\sqrt{1-2(\sin  \alpha \cos \alpha)^2}=\sqrt{1-2\cdot\left(\dfrac{2}{5}\right)^2}\\
			&=&\dfrac{\sqrt{17}}{5}\approx  0{,}82.
		\end{eqnarray*}
	} 
\end{ex}

% nguồn: Ok_TraLoiNgan_Lop11_Lan2
\begin{ex}%[1D1V3-5]%[ma:ac2d2f1b]
	Tính $S=\sin ^25^{\circ}+\sin ^210^{\circ}+\sin ^215^{\circ}+\cdots+\sin ^280^{\circ}+\sin ^285^{\circ}$. (Kết quả là số thập  phân).
	\shortans{8,5} 
	 \loigiai{
		Ta  có
		 \[
		\begin{aligned}
			& \sin^2 5^{\circ}+\sin^2  85^{\circ}=\cos^2 85^{\circ}+\sin^2 85^{\circ}=1, \\
			 & \sin^2 10^{\circ}+\sin^2 80^{\circ}=\cos^2 80^{\circ}+\sin^2 80^{\circ}=1, \\
			&  \cdots \\
			&  \sin^2 40^{\circ}+\sin^2 45^{\circ}=\cos^2 45^{\circ}+\sin^2 45^{\circ}=1.
		\end{aligned}
		\] 
		Tổng số có  $8$ cặp dư ra $\sin ^245^{\circ}$ nên $S=8+\dfrac{1}{2}=\dfrac{17}{2}$.
	 }
\end{ex}

% nguồn: Ok_TraLoiNgan_Lop11_Lan2
\begin{ex}%[1D1V3-5]%[ma:220a26f0]
	 Cho góc $\alpha $ thỏa mãn $\sin 2\alpha=-\dfrac{4}{5}$ và $\dfrac{3\pi}{4}<\alpha <\pi $. Tính $ P=\sin\alpha-\cos\alpha$ (làm tròn đến hàng phần trăm).
	 \shortans{1,34}
	\loigiai{ 
	Vì $\dfrac{3\pi}{4}<\alpha <\pi $ suy ra $\left\{\begin{array}{l}
	\sin\alpha  > 0\\
	\cos\alpha <  0
\end{array}\right.$ nên $\sin\alpha-\cos\alpha > 0$.\\
Ta có  $\left(\sin\alpha-\cos\alpha\right)^2=1-\sin 2\alpha=1+\dfrac{4}{5}=\dfrac{9}{5}$.\\
Suy ra  $\sin\alpha-\cos\alpha=\pm\dfrac{3}{\sqrt 5}$.\\
Do  $\sin\alpha-\cos\alpha > 0$ nên $\sin\alpha-\cos\alpha=\dfrac{3}{\sqrt 5}$.\\
Vậy $  P=\dfrac{3}{\sqrt 5}\approx 1{,}34$.
 }
\end{ex}

% nguồn: Ok_TraLoiNgan_Lop11_Lan2
\begin{ex}%[1D1V3-5]%[ma:ae697eb7]
	Cho góc $\alpha $ thỏa mãn $0 <\alpha <\dfrac{\pi}{2}$ và $\sin\alpha=\dfrac{2}{3}$ . Tính $P=\dfrac{1+\sin 2\alpha+\cos 2\alpha}{\sin\alpha+\cos\alpha}$  (làm tròn đến hàng phần chục).
	 \shortans{1,5}
	 \loigiai{
	Ta có $P=\dfrac{2\sin\alpha\cos\alpha+2\cos^2\alpha}{\sin\alpha+\cos\alpha}=\dfrac{2\cos\alpha\left(\sin\alpha+\cos\alpha\right)}{\sin\alpha+\cos\alpha}=2\cos\alpha $.\\ 
	Từ hệ thức  $\sin ^2\alpha+\cos ^2\alpha=1$ , suy ra $\cos\alpha=\pm\sqrt{1-\sin^2\alpha}=\pm\dfrac{\sqrt 5}{3}$.\\
	Do $0 <\alpha <\dfrac{\pi}{2}$ nên ta  chọn $\cos{\alpha}=\dfrac{\sqrt{5}}{3}\Rightarrow P=\dfrac{2\sqrt{5}}{3}\approx 1{,}5$.
	} 
\end{ex}

% nguồn: Ok_TraLoiNgan_Lop11_Lan2
\begin{ex}%[1D1V3-5]%[ma:c5f060c1]
	Tính $\cos x$ biết $\sin \left(x-\dfrac{\pi}{2}\right)+\sin  \dfrac{\pi}{2}=\sin \left(x+\dfrac{\pi}{2}\right)$.
	\shortans{0,5} 
	 \loigiai{
		Ta có \allowdisplaybreaks 
		\begin{eqnarray*}
			&& \sin \left(x-\dfrac{\pi}{2}\right)+\sin \dfrac{\pi}{2}=\sin \left(x+\dfrac{\pi}{2}\right) \\ 
			&\Leftrightarrow& -\cos  \left(\dfrac{\pi}{2}+x-\dfrac{\pi}{2}\right)+1=\cos x\\
			&\Leftrightarrow& -\cos x+1=\cos x \\ 
			&\Leftrightarrow& \cos x=\dfrac{1}{2}. 
		\end{eqnarray*}
	} 
\end{ex}

% nguồn: Ok_TraLoiNgan_Lop11_Lan2
\begin{ex}%[1D1V3-5]%[ma:9e09fcc2]
	 Cho biểu thức $D=\left(\sin ^4x+\cos ^4x-1\right)\left(\tan ^2x+\cot ^2x+2\right)$ khi $x=\dfrac{2024\pi}{2023}$ thì $D$ bằng bao nhiêu?
	 \shortans{-2}
	\loigiai{ 
		Ta  có
		 \allowdisplaybreaks
		\begin{eqnarray*}
			&&\left(\sin ^4x+\cos ^4x-1\right)\left(\tan  ^2x+\cot ^2x+2\right)\\
			&=&\left(1-2\sin ^2x \cdot \cos ^2x-1\right)\left(\dfrac{\sin ^2x}{\cos ^2x}+\dfrac{\cos  ^2x}{\sin ^2x}+2\right)\\
			&=&\left(-2\sin ^2x \cdot \cos ^2x\right)\left(\dfrac{\sin ^4x+\cos  ^4x+2\sin ^2x \cdot \cos ^2x}{\sin ^2x \cos ^2x}\right)\\
			&=&(-2)\left(\sin ^2x+\cos  ^2x\right)^2=-2.
		\end{eqnarray*}
	 }
\end{ex}

% nguồn: Ok_TraLoiNgan_Lop11_Lan2
\begin{ex}%[1D1V3-5]%[ma:b04beafc]
	Tính giá trị của  biểu thức $M=\cos^6{15^{\circ}}-\sin^6{15^{\circ}}$ (làm tròn đến hàng phần trăm).
	 \shortans{0,81}
	\loigiai{ 
	Ta  có
	$$\begin{array}{l}
	 \cos ^6\alpha-\sin ^6\alpha\\
	=\left(\cos^2\alpha-\sin^2\alpha\right)\left(\cos^4\alpha+\cos^2\alpha  \cdot\sin^2\alpha+\sin^4\alpha\right)\\
	=\cos  2\alpha \cdot\left[\left(\cos^2\alpha+\sin^2\alpha\right)^2-\cos^2\alpha \cdot\sin^2\alpha\right]\\
	 =\cos 2\alpha \cdot\left(1-\dfrac{1}{4}{\sin^2}2\alpha\right).
\end{array}$$
Vậy $ M=\cos{30^{\circ}}\cdot\left(1-\dfrac{1}{4}{\sin^2}{30^{\circ}}\right)=\dfrac{\sqrt 3}{2}\cdot\left(1-\dfrac{1}{4}\cdot\dfrac{1}{4}\right)=\dfrac{15\sqrt 3}{32}\approx  0{,}81$.
} 
\end{ex}

% nguồn: Ok_TraLoiNgan_Lop11_Lan2
\begin{ex}%[1D1V3-5]%[ma:fc3c4cfe]
	Cho biểu  thức $A=\dfrac{\left(1-\tan ^2x\right)^2}{4\tan ^2x}-\dfrac{1}{4\sin ^2x \cos ^2x}$ khi $x=\dfrac{2024\pi}{2023}$ thì $A$ bằng bao nhiêu?
	\shortans{-1} 
	 \loigiai{
		 Ta có
		\allowdisplaybreaks 
		\begin{eqnarray*}
			A&=&\dfrac{\left(1-\dfrac{\sin^2 x}{\cos^2 x}\right)^2}{4  \tan^2 x}-\dfrac{1}{4 \sin^2 x \cos^2 x}=\dfrac{\left(\cos^2 x-\sin^2 x\right)^2}{4 \sin^2 x \cos^2 x}-\dfrac{1}{4 \sin^2 x \cos^2 x} \\
			& =&\dfrac{\left(\cos^2 x-\sin^2 x+1\right)\left(\cos^2 x-\sin^2 x-1\right)}{4 \sin^2 x \cos^2 x}=\dfrac{2 \cos^2 x \cdot\left(-2 \sin^2 x\right)}{4 \sin^2  x \cos^2 x}=-1.
		\end{eqnarray*}
	 }
\end{ex}

% nguồn: Ok_TraLoiNgan_Lop11_Lan2
\begin{ex}%[1D1V3-5]%[ma:b29da4f3]
	Tính giá trị của biểu thức $ A=\sin\dfrac{\pi}{48}\cdot\cos\dfrac{\pi}{48}\cdot\cos\dfrac{\pi}{24}\cdot\cos\dfrac{\pi}{12}\cdot\cos\dfrac{\pi}{6}$ (làm tròn đến hàng  phần trăm).
	 \shortans{0,05}
	 \loigiai{
	Áp dụng công thức $\sin 2a=2\cdot\sin a\cdot\cos a,$  ta có
	\begin{align*}
		 A=&\,\sin\dfrac{\pi}{48}\cdot\cos\dfrac{\pi}{48}\cdot\cos\dfrac{\pi}{24}\cdot\cos\dfrac{\pi}{12}\cdot\cos\dfrac{\pi}{6}\\
		=&\,\dfrac{1}{2}\cdot\sin\dfrac{\pi}{24}\cdot\cos\dfrac{\pi}{24}\cdot\cos\dfrac{\pi}{12}\cdot\cos\dfrac{\pi}{6}\\  =&\,\dfrac{1}{4}\cdot\sin\dfrac{\pi}{12}\cdot\cos\dfrac{\pi}{12}\cdot\cos\dfrac{\pi}{6}\\
		 =&\,\dfrac{1}{8}\cdot\sin\dfrac{\pi}{6}\cdot\cos\dfrac{\pi}{6}\\
		 =&\,\dfrac{1}{16}\cdot\sin\dfrac{\pi}{3}=\dfrac{\sqrt 3}{32}\approx 0{,}05.
	\end{align*}
	 }
\end{ex}

% nguồn: Ok_TraLoiNgan_Lop11_Lan2
\begin{ex}%[1D1V3-5]%[ma:12314983]
	Tính giá trị của biểu  thức $P=\dfrac{\sin\dfrac{5\pi}{18}\cos\dfrac{\pi}{9}-\sin\dfrac{\pi}{9}\cos\dfrac{5\pi}{18}}{\cos\dfrac{\pi}{4}\cos\dfrac{\pi}{12}-\sin\dfrac{\pi}{4}\sin\dfrac{\pi}{12}}$.
	\shortans{1} 
	 \loigiai{
	Áp dụng công thức $\left\{\begin{array}{l}
	\sin  a\cdot\cos b-\cos a\cdot\sin b=\sin\left(a-b\right)\\
	\cos a\cdot\cos  b-\sin a\cdot\sin b=\cos\left(a+b\right)
\end{array}\right.$\\
 Khi đó $\sin\dfrac{5\pi}{18}\cos\dfrac{\pi}{9}-\sin\dfrac{\pi}{9}\cos\dfrac{5\pi}{18}=\sin\left(\dfrac{5\pi}{18}-\dfrac{\pi}{9}\right)=\sin\dfrac{\pi}{6}=\dfrac{1}{2}.$\\
Và $\cos\dfrac{\pi}{4}\cos\dfrac{\pi}{12}-\sin\dfrac{\pi}{4}\sin\dfrac{\pi}{12}=\cos\left(\dfrac{\pi}{4}+\dfrac{\pi}{12}\right)=\cos\dfrac{\pi}{3}=\dfrac{1}{2}$.\\ 
Vậy  $P=\dfrac{1}{2}:\dfrac{1}{2}=1$.
} 
\end{ex}

% nguồn: Ok_TraLoiNgan_Lop11_Lan2
\begin{ex}%[1D1V3-5]%[ma:2678670d]
	Cho biểu thức $C=\cos ^2x+\cos  ^2\left(\dfrac{\pi}{3}+x\right)+\cos ^2\left(\dfrac{\pi}{3}-x\right)$ khi $x=\dfrac{2024\pi}{2023}$ thì $C$ bằng bao nhiêu?
	\shortans{1,5} 
	 \loigiai{
		 Ta có
		\[ 
		\begin{aligned}
			C  &=\cos^2 x+\dfrac{1+\cos \left(\dfrac{2\pi}{3}-2 x\right)}{2}+\dfrac{1+\cos \left(\dfrac{2\pi}{3}+2 x\right)}{2} \\
			&=\cos^2 x+1+\cos \left(\dfrac{2\pi}{3}\right) \cos  (-2 x)=\dfrac{1+\cos 2 x}{2}+1+\left(-\dfrac{1}{2}\right) \cos 2 x\\
			&=\dfrac{3}{2}=1{,}5. 
		\end{aligned}
		 \]
	} 
\end{ex}

% nguồn: Ok_TraLoiNgan_Lop11_Lan2
\begin{ex}%[1D1V3-5]%[ma:bacec3ce]
	Tính giá trị đúng của biểu thức $\dfrac{\tan{225^{\circ}}-\cot{81^{\circ}}\cot{69^\circ}}{\cot{261^{\circ}}+\tan{201^{\circ}}}$ (làm tròn đến hàng phần  trăm).
	 \shortans{1,73}
	 \loigiai{
	Ta  có
	\begin{align*}
		\dfrac{\tan{225^{\circ}}-\cot{81^{\circ}}.cot{69^{\circ}}}{\cot{261^{\circ}}+\tan{201^{\circ}}}=&\,\dfrac{\tan\left(180^{\circ}+45^{\circ}\right)-\tan{9^{\circ}}\cdot\cot{69^{\circ}}}{\cot\left(180^{\circ}+81^{\circ}\right)+\tan\left(180^{\circ}+21^{\circ}\right)}\\ 
		 =&\,\dfrac{1-\tan{9^{\circ}}\cdot\tan{21^{\circ}}}{\tan{9^{\circ}}+\tan{21^{\circ}}}=\dfrac{1}{\tan\left(9^{\circ}+21^{\circ}\right)}\\
		=&\,\dfrac{1}{\tan{30^{\circ}}}=\sqrt  3.
	\end{align*}
	 }
\end{ex}

% nguồn: Ok_TraLoiNgan_Lop11_Lan2
\begin{ex}%[1D1V3-5]%[ma:7d01b9dc]
	Cho $3\cos \alpha-\sin \alpha=1$, với  $0^{\circ} < \alpha < 90^{\circ}$. Tính giá trị của $\tan \alpha$. (Kết quả làm tròn đến hàng phần mười).
	 \shortans{1,3}
	 \loigiai{
		 Ta có \allowdisplaybreaks
		\begin{eqnarray*}
			&&3\cos \alpha-\sin  \alpha=1\\
			&\Leftrightarrow& 3\cos  \alpha=\sin \alpha+1\\
			&\Rightarrow&  9\cos ^2\alpha=(\sin \alpha+1)^2\\
			 &\Leftrightarrow& 9\cos ^2\alpha=\sin ^2\alpha+2\sin \alpha+1\\
			&\Leftrightarrow& 9\left(1-\sin ^2\alpha\right)=\sin ^2\alpha+2\sin  \alpha+1\\
			&\Leftrightarrow&  10\sin ^2\alpha+2\sin \alpha-8=0\\
			&\Leftrightarrow& \hoac{&\sin  \alpha=-1\\
			&\sin  \alpha=\dfrac{4}{5}.}
		\end{eqnarray*}
		\begin{itemize}
			\item $\sin \alpha=-1$: không thỏa mãn vì $0^{\circ} < \alpha < 90^{\circ}$. 
			\item $\sin \alpha=\dfrac{4}{5} \Rightarrow \cos \alpha=\dfrac{3}{5} \Rightarrow \tan \alpha=\dfrac{\sin \alpha}{\cos \alpha}=\dfrac{4}{3}\approx 1{,}3$. 
		\end{itemize}
	 }
\end{ex}

% nguồn: Ok_TraLoiNgan_Lop11_Lan2
\begin{ex}%[1D1V3-5]%[ma:8e175466]
	Tính $\sin^2  x$ biết $\cot \left(x+540^{\circ}\right)-\tan \left(x-90^{\circ}\right)=\sin ^2\left(725^{\circ}\right)+\cos ^2\left(365^{\circ}\right)$.
	 \shortans{0,8}
	\loigiai{ 
		Ta có 
		 \allowdisplaybreaks
		\begin{eqnarray*}
			&&  \cot \left(x+540^{\circ}\right)-\tan \left(x-90^{\circ}\right)=\sin^2\left(725^{\circ}\right)+\cos^2\left(365^{\circ}\right) \\
			& \Leftrightarrow&  \cot \left(x+180^{\circ}+360^{\circ}\right)-\tan \left(x-90^{\circ}\right)=\sin^2\left(725^{\circ}+(-1)\cdot 360^{\circ}\right)+\cos^2\left(365^{\circ}\right) \\
			&  \Leftrightarrow &\cot \left(x+180^{\circ}\right)+\cot \left(90^{\circ}+x-90^{\circ}\right)=\sin^2\left(365^{\circ}\right)+\cos^2\left(365^{\circ}\right) \\
			& \Leftrightarrow &\cot x+\cot  x=1 \Leftrightarrow \cot x=\dfrac{1}{2}.
		\end{eqnarray*}
		Khi đó $\sin^2  x=\dfrac{1}{1+\cot^2 x}=\dfrac{4}{5}=0{,}8$.
	} 
\end{ex}

% nguồn: Ok_TraLoiNgan_Lop11_Lan2
\begin{ex}%[1D1V3-5]%[ma:b5902086]
	Tính giá trị của biểu thức $M=\cos^4{15^\circ}-\sin ^4{15^\circ}+\cos^2{15^\circ}-\sin^2{15^\circ}$ (làm trong  đến hàng phần trăm).
	 \shortans{1,73}
	 \loigiai{
	Áp dụng công thức nhân đôi  $\cos ^2a-\sin ^2a=\cos 2a$ .\\
	Ta  có
	\begin{align*}
		M&\,=\left(\cos^415^{\circ}-\sin^415^{\circ}\right)+\left(\cos^215^{\circ}-\sin^215^{\circ}\right)\\ 
		 &\,=\left(\cos^215^{\circ}-\sin^215^{\circ}\right)\left(\cos^215^{\circ}+\sin^215^{\circ}\right)+\left(\cos^215^{\circ}-\sin^215^{\circ}\right)\\
		&\,=\left(\cos^215^{\circ}-\sin^215^{\circ}\right)+\left(\cos^215^{\circ}-\sin^215^{\circ}\right)\\ 
		 &\,=\cos{30^{\circ}}+\cos{30^{\circ}}=\sqrt 3.
	\end{align*}
	} 
\end{ex}

% nguồn: Ok_TraLoiNgan_Lop11_Lan2
\begin{ex}%[1D1V3-5]%[ma:21372554]
	 Cho biểu thức $B=\left[\sin \left(\dfrac{\pi}{2}-x\right)+\sin (10\pi+x)\right]^2+\left[\cos \left(\dfrac{3\pi}{2}-x\right)+\cos (8\pi-x)\right]^2$ khi $x=\dfrac{2024\pi}{2023}$ thì $B$ bằng bao nhiêu?
	 \shortans{2}
	 \loigiai{
		Ta  có
		 \allowdisplaybreaks
		\begin{eqnarray*}
			\sin \left(\dfrac{\pi}{2}-x\right)=\cos x,\,  \sin (10\pi+x)=\sin x,\, \cos \left(\dfrac{3\pi}{2}-x\right)=-\sin x,\, \cos (8\pi-x)=\cos x.
		\end{eqnarray*}
		Khi đó \allowdisplaybreaks 
		\begin{eqnarray*}
			B&=&\left[\sin \left(\dfrac{\pi}{2}-x\right)+\sin (10\pi+x)\right]^2+\left[\cos \left(\dfrac{3\pi}{2}-x\right)+\cos  (8\pi-x)\right]^2\\
			&=&(\cos x+\sin x)^2+(-\sin  x+\cos x)^2\\
			&=&2. 
		\end{eqnarray*}
	} 
\end{ex}

% nguồn: Ok_TraLoiNgan_Lop11_Lan2
\begin{ex}%[1D1V3-5]%[ma:8e30aed6]
	Cho góc $\alpha  $ thỏa mãn $\sin\alpha=\dfrac{3}{5}$. Tính $P=\sin\left(\alpha+\dfrac{\pi}{6}\right)\sin\left(\alpha-\dfrac{\pi}{6}\right)$ (làm tròn đến hàng phần trăm).
	 \shortans{0,11}
	 \loigiai{
	Áp dụng công thức $\sin a\cdot\sin  b=\dfrac{1}{2}\left[\cos\left(a-b\right)-\cos\left(a+b\right)\right]$, ta được
	$$P=\sin\left(\alpha+\dfrac{\pi}{6}\right)\sin\left(\alpha-\dfrac{\pi}{6}\right)=\dfrac{1}{2}\left(\cos\dfrac{\pi}{3}-\cos  2\alpha\right).$$
	Ta  có $\cos 2\alpha=1-2\sin ^2\alpha=1-2\cdot\left(\dfrac{3}{5}\right)^2=\dfrac{7}{25}.$\\
	Thay vào $ P$, ta được $P=\dfrac{1}{2}\left(\dfrac{1}{2}-\dfrac{7}{25}\right)=\dfrac{11}{100}\approx 0{,}11$. 
	} 
\end{ex}

% nguồn: Ok_TraLoiNgan_Lop11_Lan2
\begin{ex}%[1D1V3-5]%[ma:9ca070a5]
	Cho góc $\alpha $ thỏa mãn $\sin\alpha=\dfrac{4}{5}$. Tính $P=\cos 4\alpha$ (làm tròn đến hàng phần  chục).
	\shortans{-0,8} 
	 \loigiai{
	Ta có $\cos  {2\alpha}=1-2\sin^2\alpha=1-2\cdot\left(\dfrac{4}{5}\right)^2=-\dfrac{7}{25}.$\\
	Suy ra $  P=\cos 4\alpha=2\cos ^22\alpha-1=2\cdot\dfrac{49}{625}-1=-\dfrac{527}{625}=-0{,}8$.
	} 
\end{ex}


%%%%%%%%%%%%%%%% TỰ LUẬN %%%%%%%%%%%%%%%%

% nguồn: Ok_TuanLuan_Toan11_Lan2
\begin{ex}%[1D1V3-5]%[ma:17cb82b4]
	Chứng minh các đẳng  thức sau
	\begin{enumerate}
		\item $\sin ^4a+\cos  ^4a=1-\dfrac{1}{2}{\sin ^2}2a=\dfrac{3}{4}+\dfrac{1}{4}\cos 4a$;
		 \item $\cos a-\sin a=\sqrt 2\cos\left(a+\dfrac{\pi}{4}\right)$;
		\item  $\sin a+\sqrt 3\cos a=2\sin\left(a+\dfrac{\pi}{3}\right)$.
	\end{enumerate}
	\loigiai{ 
	\begin{enumerate}
		\item Ta có $$\begin{array}{l}
		 \sin ^4a+\cos ^4a\\
		 =\left(\sin^2a+\cos^2a\right)^2-2\sin ^2a{\cos ^2}a\\
		=1-2\cdot\left(\sin  a\cdot\cos a\right)^2 \\
		 =1-\dfrac{1}{2}{\sin ^2}2a=1-\dfrac{1}{2}\left(\dfrac{1-\cos4a}{2}\right)\\
		 =1-\dfrac{1}{4}+\dfrac{1}{4}\cos 4a=\dfrac{3}{4}+\dfrac{1}{4}\cos 4a.
	\end{array}$$
	\item Ta có $$\begin{array}{l}
	 \sqrt 2\cos\left(a+\dfrac{\pi}{4}\right)\\
	=\sqrt 2\left(\cos a\cos\dfrac{\pi}{4}-\sin a\sin\dfrac{\pi}{4}\right)\\ 
	=\sqrt 2\left(\dfrac{\sqrt 2}{2}\cos a-\dfrac{\sqrt 2}{2}\sin a\right)\\ 
	=\sqrt 2 \cdot\dfrac{\sqrt  2}{2}\cdot\left(\cos a-\sin a\right)=\cos a-\sin a.
\end{array}$$
\item Ta  có\\
$2\sin\left(a+\dfrac{\pi}{3}\right)=2\left(\sin a\cos\dfrac{\pi}{3}+\cos a\sin\dfrac{\pi}{3}\right)=2\left(\dfrac{1}{2}\sin a+\dfrac{\sqrt  3}{2}\cos a\right)=\sin a+\sqrt 3\cos a{\rm}$
\end{enumerate}
 }
\end{ex}

% nguồn: Ok_TuanLuan_Toan11_Lan2
\begin{ex}%[1D1V3-5]%[ma:6778351b]
	Chứng  minh các đẳng thức lượng giác sau
	\begin{multicols}{2}
		\begin{enumerate}
			\item  $\sin\left(60^{\circ}+\alpha\right)-\sin{\left(60^{\circ}-\alpha\right)}=\sin{\alpha}$;
			\item $\sin ^4\alpha+\cos ^4\alpha=\dfrac{3}{4}+\dfrac{1}{4}\cos  4\alpha $;
			\item $\sin\alpha\left(2\cos 4\alpha+2\cos 2\alpha+1\right)=\sin 5\alpha $; 
			\item  $\dfrac{\cos\left(\alpha-\beta\right)}{\cos\left(\alpha+\beta\right)}=\dfrac{1+\tan\alpha\tan\beta}{1-\tan\alpha\tan\beta}$.
		\end{enumerate}
	\end{multicols}
	 \loigiai{
	\begin{enumerate}
		\item Ta có $$\begin{array}{l}
		\sin\left(60^\circ+\alpha\right)-\sin\left(60^\circ-\alpha\right)\\ 
		=\left(\sin  60^\circ\cos\alpha+\cos 60^\circ\sin\alpha\right)-\left(\sin 60^\circ\cos\alpha-\cos 60^\circ\sin\alpha\right)\\
		=2\cos60^\circ\sin\alpha\\ 
		 =2\cdot\dfrac{1}{2}\sin\alpha=\sin\alpha.
	\end{array}$$
	\item Ta  có
	$$\begin{array}{l}
	\sin ^4\alpha+\cos  ^4\alpha\\
	 =\left(\sin^2\alpha+\cos^2\alpha\right)^2-2\sin ^2\alpha{\cos ^2}\alpha \\
	=1-\dfrac{1}{2}{\sin ^2}2\alpha\\ 
	=1-\dfrac{1}{4}\left(1-\cos 4\alpha\right)\\ 
	 =1-\dfrac{1}{4}+\dfrac{1}{4}\cos 4\alpha\\
	=\dfrac{3}{4}+\dfrac{1}{4}\cos 4\alpha. 
\end{array}$$
\item Ta có $$\begin{array}{l}
\sin\alpha\left(2\cos 4\alpha+2\cos  2\alpha+1\right)\\
 =2\sin\alpha\cos 4\alpha+2\sin\alpha\cos 2\alpha+\sin\alpha\\
=2\cdot\dfrac{1}{2}\left[\sin\left(-3\alpha\right)+\sin 5\alpha\right]+2\cdot\dfrac{1}{2}\left[\sin\left(-\alpha\right)+\sin 3\alpha\right]+\sin\alpha\\ 
=\left(-\sin 3\alpha+\sin 5\alpha\right)+\left(-\sin\alpha+\sin  3\alpha\right)+\sin\alpha\\
 =\sin 5\alpha.
\end{array}$$
\item Ta có $$\begin{array}{l}
\dfrac{\cos\left(\alpha-\beta\right)}{\cos\left(\alpha+\beta\right)}=\dfrac{\cos\alpha\cos\beta+\sin\alpha\sin\beta}{\cos\alpha\cos\beta-\sin\alpha\sin\beta} 
 =\dfrac{\cos\alpha\cos\beta\left(1+\tan\alpha\tan\beta\right)}{\cos\alpha\cos\beta\left(1-\tan\alpha\tan\beta\right)}
=\dfrac{1+\tan\alpha\tan\beta}{1-\tan\alpha\tan\beta}. 
\end{array}$$
\end{enumerate}
 }
\end{ex}

% nguồn: Ok_TuanLuan_Toan11_Lan2
\begin{ex}%[1D1V3-5]%[ma:6aed4e37]
	Cho tam giác $ ABC$, chứng minh  rằng
	\begin{enumerate}
		\item $\tan A+\tan B+\tan C=\tan A\cdot\tan B\cdot\tan C$ (với  điều kiện tam giác $ ABC$ không vuông);
		 \item $\tan\dfrac{A}{2}\cdot\tan\dfrac{B}{2}+\tan\dfrac{B}{2}\cdot\tan\dfrac{C}{2}+\tan\dfrac{C}{2}\cdot\tan\dfrac{A}{2}=1$.
	\end{enumerate}
	\loigiai{ 
	\begin{enumerate}
		\item Do $A+B=\pi-C$ nên 
		\begin{align*}
			&\,\tan{(A+B)}=-\tan{C}\Leftrightarrow\dfrac{\tan  A+\tan B}{1-\tan A\tan B}=-\tan C\\
			\Leftrightarrow&\,\tan A+\tan B+\tan C=\tan A\cdot\tan  B\cdot\tan C.
		\end{align*}
		\item Do  $\dfrac{A}{2}+\dfrac{B}{2}=\dfrac{\pi}{2}-\dfrac{C}{2}$ nên
		\begin{align*}
			 &\,\tan{\left(\dfrac{A}{2}+\dfrac{B}{2}\right)}=\cot{\dfrac{C}{2}}\Leftrightarrow\dfrac{\tan\dfrac{A}{2}+\tan\dfrac{B}{2}}{1-\tan\dfrac{A}{2}\tan\dfrac{B}{2}}=\dfrac{1}{\tan\dfrac{C}{2}}\\
			 \Leftrightarrow&\,\tan\dfrac{A}{2}\cdot\tan\dfrac{B}{2}+\tan\dfrac{B}{2}\cdot\tan\dfrac{C}{2}+\tan\dfrac{C}{2}\cdot\tan\dfrac{A}{2}=1.
		\end{align*}
	\end{enumerate}
	 }
\end{ex}

% nguồn: Ok_TuanLuan_Toan11_Lan2
\begin{ex}%[1D1V3-5]%[ma:5ac482c2]
	Rút gọn  các biểu thức sau
	\begin{multicols}{3}
		\begin{enumerate}
			\item $\dfrac{\sin\alpha+\sin  2\alpha}{1+\cos\alpha+\cos 2\alpha}$;
			 \item $\dfrac{\cos\left(\dfrac{\pi}{4}+\alpha\right)-\cos\left(\dfrac{\pi}{4}-\alpha\right)}{\sin\left(\dfrac{\pi}{4}+\alpha\right)-\sin\left(\dfrac{\pi}{4}-\alpha\right)}$;
			\item $\dfrac{\sin^2\alpha}{4-4\sin^2\dfrac{\alpha}{2}}$. 
		\end{enumerate}
	\end{multicols}
	 \loigiai{
	\begin{enumerate}
		\item Ta có $\dfrac{\sin\alpha+\sin 2\alpha}{1+\cos\alpha+\cos 2\alpha}=\dfrac{\sin\alpha+2\sin\alpha\cos\alpha}{1+\cos\alpha+2\cos^2\alpha-1}=\dfrac{\sin\alpha  (1+2\cos\alpha)}{\cos\alpha (1+2\cos\alpha)}=\dfrac{\sin\alpha}{\cos\alpha}=\tan\alpha $;
		\item  Ta có $\dfrac{\cos\left(\dfrac{\pi}{4}+\alpha\right)-\cos\left(\dfrac{\pi}{4}-\alpha\right)}{\sin\left(\dfrac{\pi}{4}+\alpha\right)-\sin\left(\dfrac{\pi}{4}-\alpha\right)}=\dfrac{-2\sin\dfrac{\dfrac{\pi}{2}}{2}\sin\dfrac{2\alpha}{2}}{2\cos\dfrac{\dfrac{\pi}{2}}{2}\sin\dfrac{2\alpha}{2}}=-\dfrac{\sin\dfrac{\pi}{4}\sin\alpha}{\cos\dfrac{\pi}{4}\sin\alpha}=-\tan\dfrac{\pi}{4}=-1$;
		\item Ta có $\dfrac{\sin^2\alpha}{4-4\sin^2\dfrac{\alpha}{2}}=\dfrac{\left(2\sin\dfrac{\alpha}{2}\cos\dfrac{\alpha}{2}\right)^2}{4\left(1-\sin^2\dfrac{\alpha}{2}\right)}=\dfrac{4\sin^2\dfrac{\alpha}{2}{\cos^2}\dfrac{\alpha}{2}}{4\cos^2\dfrac{\alpha}{2}}=\sin ^2\dfrac{\alpha}{2}$. 
	\end{enumerate}
	 }
\end{ex}

% nguồn: Ok_TuanLuan_Toan11_Lan2
\begin{ex}%[1D1V3-5]%[ma:471b7246]
	Cho $\cos\left(a+2b\right)=2\cos a$. Chứng minh rằng $\tan\left(a+b\right)\cdot\tan  b=-\dfrac{1}{3}$.
	\loigiai{ 
	 Ta có
	\begin{align*}
		&\,\cos\left(a+2b\right)=2\cos  a\\
		 \Leftrightarrow&\,\cos\left[\left(a+b\right)+b\right]=2\cos\left[\left(a+b\right)-b\right]\\
		 \Leftrightarrow&\, 3\sin\left(a+b\right)\cdot\sin b=-\cos\left(a+b\right)\cdot\cos b\\
		 \Leftrightarrow&\,\tan\left(a+b\right)\cdot\tan b=-\dfrac{1}{3}.
	\end{align*}
	Vậy đẳng thức được  chứng minh.
	} 
\end{ex}
